\documentclass[10pt,a4paper]{article}

% ----------------- Paquetes -------------------%
\usepackage[T1]{fontenc}
\usepackage[utf8]{inputenc}
\usepackage[a4paper,margin=2.5cm]{geometry}
\usepackage{mathptmx}             
\usepackage{changepage} 
\usepackage{setspace}
\usepackage[spanish]{babel}
\usepackage[labelfont=bf,labelsep=space]{caption}
\usepackage{amssymb}
\usepackage{amsmath}
\usepackage{ebgaramond}
\usepackage{placeins}
\usepackage{etoolbox}
\usepackage{setspace}
\usepackage{parskip} 
\usepackage{tcolorbox}
\usepackage{needspace}
\usepackage{xparse}
\usepackage{ocgx2}
\usepackage{hyperref}
\usepackage{hypcap}
\usepackage{soul}\sethlcolor{yellow}% resaltado amarillo: \hl{texto} (Panel TCEE, ⌃H)


%%%%%%%%%%%%%%%%%%%%%%%%%%%%%%%%%%%%%%%%%%%%%%%%%%%%%%%%%%%%%%%%
% --- Fija primero tus valores globales "buenos" ---
\setstretch{1.2}                 % Interlineado global
\setlength{\parskip}{0.7\baselineskip}
\setlength{\parindent}{0pt}
\newlength{\GlobalParskip}
\setlength{\GlobalParskip}{\parskip}

\newcounter{eqnum}[subsection]
\renewcommand{\theeqnum}{\arabic{eqnum}}
\newcounter{alphaitem}
\newcounter{numitem}

%% ------------------- Recuadros----------------------%%
\newcommand{\puntosclave}[1]{%
  \begin{tcolorbox}[
    colframe=black,
    title={Puntos clave},
    before upper={\setlength{\parindent}{0.5cm}\setlength{\parskip}{0.5\baselineskip}}
  ]
  #1
  \end{tcolorbox}
}

\newcommand{\modificaciones}[1]{
  \begin{tcolorbox}[
    colback=white,
    colframe=red,
    title={Modificaciones a realizar},
    before upper={\setlength{\parindent}{0.5cm}\setlength{\parskip}{0.5\baselineskip}}
  ]
  #1
  \end{tcolorbox}
}

%% ------------------- Preguntas TEST----------------------%%
% Opciones: radio (single-choice)
\NewDocumentCommand{\OptRadio}{m m m}{%
  \noindent\CheckBox[radio,name=#1,value=#2,width=1.2ex,height=1.2ex]{}%
  \;\textbf{#2.}~#3\par
}

% Opciones: checkbox (multi-choice)
\NewDocumentCommand{\OptCheck}{m m}{%
  \noindent\CheckBox[name=#1,width=1.2ex,height=1.2ex,bordercolor={0 0 0}]{}%
  \;#2\par
}

% Pregunta single-choice (radio)
% Args: 1=id, 2=enunciado, 3=opciones, 4=solución+justificación, 5=imagen (o {}), 6=origen
\NewDocumentCommand{\PreguntaTestSingle}{m m m m m m}{%
  \par\medskip
  \noindent\textbf{#1.}~#2\par\smallskip

    % Imagen (si no es {})
    \IfBlankTF{#5}{}{%
      \begin{center}
        \includegraphics[width=0.75\linewidth]{#5}
      \end{center}
    }

    \begin{Form}
      #3
    \end{Form}

    \par\smallskip
    \switchocg{sol-#1}{\fbox{Mostrar / ocultar solución}}%
    \hfill{\footnotesize #6}\par\smallskip

    \begin{ocg}{Solución}{sol-#1}{0}
      \noindent #4
    \end{ocg}
  \par\medskip
}

% Pregunta multi-choice (checkboxes)
\NewDocumentCommand{\PreguntaTestMulti}{m m m m m m}{%
  \par\medskip
  \noindent\textbf{#1.}~#2\par\smallskip

    % Imagen (si no es {})
    \IfBlankTF{#5}{}{%
      \begin{center}
        \includegraphics[width=0.75\linewidth]{#5}
      \end{center}
    }
    \begin{Form}
      #3
    \end{Form}

    \par\smallskip
    \switchocg{sol-#1}{\fbox{Mostrar / ocultar solución}}%
    \hfill{\footnotesize #6}\par\smallskip

    \begin{ocg}{Solución}{sol-#1}{0}
      \noindent #4
    \end{ocg}
  \par\medskip
}

%% ------------------- Listas----------------------%%
    % --- Lista alfabética ---
    \newenvironment{la}
      {\begin{list}{\alph{alphaitem})}{%
          \usecounter{alphaitem}%
          \setlength{\leftmargin}{1cm}%
          \setlength{\parskip}{3pt}% 
          \setlength{\itemsep}{0pt}% ← espacio entre ítems
          \setlength{\parsep}{0pt}%  ← párrafos dentro de un ítem
      }}
      {\end{list}}
      
    % --- Lista numérica ---
    \newenvironment{lnum}
      {\begin{list}{\arabic{numitem}.}{%
          \usecounter{numitem}%
          \setlength{\leftmargin}{1cm}%
          \setlength{\parskip}{3pt}% 
          \setlength{\itemsep}{0pt}% ← espacio entre ítems
          \setlength{\parsep}{0pt}%  ← párrafos dentro de un ítem
      }}
      {\end{list}}
      
% ------------------- Imágenes----------------------%
    \graphicspath{{Img/}}% Rutas por defecto 
    % --- Imagen adaptativa ---
    % Registros de dimensión definidos una sola vez
    \newdimen\imgcap
    \newdimen\imgmin
    \newdimen\imgmax
    \newdimen\imgremain
    \newdimen\imgh
    \newdimen\imgneed
    
    \newcommand{\imagenfit}[3]{%
      \addvspace{10pt}% separación antes
      \par\begingroup
        % Parámetros de composición (asignaciones, no \newdimen)
        \imgcap=1\baselineskip            % alto estimado de título+fuente
        \imgmin=0.15\textheight
        \imgmax=0.15\textheight
        \imgremain=\pagegoal
        \ifdim\pagegoal=\maxdimen \imgremain=0pt\fi
        \advance\imgremain by -\pagetotal
        \advance\imgremain by -\imgcap
        % Decisión de altura destino
        \ifdim\imgremain<\imgmin
          \newpage
          \imgh=0.20\textheight
        \else
          \imgh=\imgremain
          \ifdim\imgh>\imgmax \imgh=\imgmax \fi
        \fi
        % Reservar espacio para evitar cortes del bloque completo
        \imgneed=\imgh \advance\imgneed by \imgcap
        \Needspace{\imgneed}
        % Cuerpo
        \noindent\begin{minipage}{\textwidth}
          \centering
          \captionof{figure}{\textemdash\ #2}\par\vspace{0.3em}% título arriba
          \includegraphics[width=\linewidth,height=\imgh,keepaspectratio]{\detokenize{#1}}\par
          \vspace{0.3em}\small\textit{Fuente: }#3% fuente abajo
        \end{minipage}%
      \endgroup
      \par\addvspace{10pt}%
    }

%------------ Bloque de RECUADRO ESPECIAL -------------%
    \newenvironment{specialblock}[1]{%
      \begin{quote} % crea sangría a izquierda y derecha
      \small % texto más pequeño
      \noindent\textbf{#1}\\[-0.3em] % título en negrita
      \rule{\linewidth}{0.4pt}\\[0.6em]}{
      \end{quote}}
      
%-------------------------- Portada -----------------------%
\newcommand{\oldtitle}[1]{\def\OldTitle{#1}}
\newcommand{\changerationale}[1]{\def\ChangeWhy{#1}}

\newcommand{\changebox}{%
\begin{tcolorbox}[title={Historial de cambios del tema}]
\textbf{Título anterior:} \OldTitle\par
\textbf{Motivo del cambio:} \ChangeWhy
\end{tcolorbox}
}

%-------------------------- Author Font -----------------------%
    \newcommand{\authorfont}[1]{{\fontfamily{EBGaramond-TLF}\selectfont #1}}

%-------------------------- Anexos -----------------------%
    \makeatletter
    \newcounter{anexo}
    \renewcommand{\theanexo}{\Roman{anexo}}
    \newcommand{\anexo}[1]{%
      \clearpage
      \phantomsection
      \refstepcounter{anexo}%
      \noindent\textbf{Anexo \theanexo.- }#1\par\medskip
      \addcontentsline{stc}{section}{Anexo \theanexo.- #1}%
    }
    \makeatother

    %-------------------------- Lista de figuras (personal) -----------------------%
\makeatletter
\newcommand{\MyListOfFigures}{%
  \clearpage
  \phantomsection
  {\centering\bfseries\Large Índice de figuras\par}%
  \medskip
  % Entrada en nuestro índice de contenido (stc)
  \addcontentsline{stc}{section}{Índice de figuras}%
  % Recuperación del listado (sin cabecera propia)
  \begingroup
    \parindent=0pt\parskip=2pt
    \@starttoc{lof}%
  \endgroup
}
\makeatother

%-------------------------- Bibliografía -----------------------%
\makeatletter
% Contador para las referencias
\newcounter{myref}
% Cabecera de la bibliografía, entra en el índice 'stc'
\newcommand{\MyBibliography}{%
  \clearpage
  \phantomsection
  {\centering\bfseries\Large Bibliografía y referencias\par}%
  \medskip
  \addcontentsline{stc}{section}{Bibliografía y referencias}%
  \setcounter{myref}{0}%
}
\newcommand{\MyBibItem}[4]{%
  \refstepcounter{myref}%
  \noindent[\themyref]\ #1.\ \emph{#2}. #3. #4\par
}
\makeatother

%--------- Establecer cuatro niveles de índice --------%
    \setcounter{tocdepth}{4}   
    \setcounter{secnumdepth}{4}% numera hasta \paragraph

%%%%%%%%%%%%%%%%%%%%%%%%%%%%%%%%

\makeatletter

    %--------- Interlineado Global --------%
    \edef\GlobalBaseLineStretch{\baselinestretch}
    
    %--------- Espaciado de seccinones --------%
    \renewcommand\section{\@startsection{section}
      {1}{\z@}
      {2.5ex \@plus 1ex \@minus .2ex} % espacio antes
      {1.5ex \@plus .2ex}             % espacio después
      {\normalfont\Large\bfseries}    % formato del título
    }
    \renewcommand\subsection{\@startsection{subsection}{2}{\z@}
      {3ex \@plus 0.8ex \@minus .2ex}
      {1.5ex \@plus .2ex}
      {\normalfont\large\bfseries}
    }
    \renewcommand\subsubsection{\@startsection{subsubsection}{3}{\z@}
      {3ex \@plus 0.5ex \@minus .2ex}
      {1ex \@plus .2ex}
      {\normalfont\normalsize\bfseries}
    }
    \renewcommand\paragraph{\@startsection{paragraph}{4}{\z@}%
    {-2ex \@plus -0.5ex \@minus -.2ex}% espacio antes
    {0.8ex \@plus .2ex}% espacio después
    {\normalfont\normalsize\itshape}}% ← cursiva en lugar de negrita

    %--------- Ecuaciones --------%   
    % Variables globales para almacenar títulos y notas
    \gdef\leq@current@section{}%
    \gdef\leq@current@subsection{}%
    \gdef\leq@last@printed@section{}%
    \gdef\leq@last@printed@subsection{}%
    
    % Comando para añadir nota (invisible en el cuerpo)
    \newcommand{\eqnote}[1]{%
      \addcontentsline{leq}{leqnote}{#1}%
    }
    
    % Comando para ecuaciones
    \newcommand{\eqblock}[2]{%
      % Verificar si necesitamos imprimir el título de sección
      \ifx\leq@current@section\leq@last@printed@section\else
        \addcontentsline{leq}{leqsection}{\leq@current@section}%
        \global\let\leq@last@printed@section\leq@current@section
        \gdef\leq@last@printed@subsection{}% Reset subsección al cambiar sección
      \fi
      % Verificar si necesitamos imprimir el título de subsección
      \ifx\leq@current@subsection\leq@last@printed@subsection\else
        \ifx\leq@current@subsection\@empty\else
          \addcontentsline{leq}{leqsubsection}{\leq@current@subsection}%
          \global\let\leq@last@printed@subsection\leq@current@subsection
        \fi
      \fi
      % Ecuación
      \refstepcounter{eqnum}%
      \begin{equation*}
        #1\tag{E.\theeqnum}
      \end{equation*}%
      \addcontentsline{leq}{leqt}{[E.\theeqnum] -- #2}%
      \addcontentsline{leq}{leqm}{#1}%
    }
    
    %%% ----- Índice de ecuaciones ----------%%%
    % Tamaño para []
    \newcommand{\leqIndexFont}{\normalsize}
    \newcommand{\leqTitleFont}{\tiny}
    \newcommand{\leqNoteFont}{\tiny}
    % Cabecera de sección 
    \newcommand*\l@leqsection[2]{%
      \addvspace{25pt}% Mayor separación antes de nueva sección
      \noindent\leqIndexFont
      \makebox[\linewidth]{\rule{.38\linewidth}{0.4pt}\ \ \textbf{  \MakeUppercase{#1}}\ \ \rule{.38\linewidth}{0.4pt}}%
      \par\vspace{-10pt}
    }
    % Cabecera de subsección 
    \newcommand*\l@leqsubsection[2]{%
      \par\addvspace{15pt}%
      \noindent\leqIndexFont #1
      \par\vspace{-7pt}
    }
    % Título de ecuación 
    \newcommand*\l@leqt[2]{%
    \par\addvspace{10pt}%
      \par\noindent\leqTitleFont #1\par
    }
    % Miniatura de ecuación
    \newcommand*\l@leqm[2]{%
      \vspace{0.3\baselineskip}%
      \par\scriptsize\noindent\hspace*{1.5em}\(#1\)\par%
    }
    % Nota de ecuación 
    \newcommand*\l@leqnote[2]{%
      \vspace{0.2\baselineskip}%
      \par\noindent\leqNoteFont\hspace*{1.5em}\textbf{Nota.--} #1\par\vspace{0.5\baselineskip}%
    }
    % Generación del índice
    \newcommand{\mylistofequations}{%
      \clearpage
      \phantomsection
      % Título visual de la lista (ajústalo a tu gusto):
      {\centering\bfseries\Large Índice de ecuaciones\par}
      % Entrada en nuestro índice 'stc':
      \addcontentsline{stc}{section}{Índice de ecuaciones}%
      % Llamada a tu lista real (usa el nombre exacto de tu macro):
      \@starttoc{leq}%
    }
    
    %--------- Captura de títulos de sección --------%
    \let\leq@original@section\section
    \let\leq@original@subsection\subsection
    % Flag para saber si estamos en una sección asterisco
    \newif\ifLeqStarSection
    \LeqStarSectionfalse
    
    % Redefinir \section CON CONTROL DE ASTERISCO
    \renewcommand{\section}{%
      \@ifstar{\leq@section@star}{\@dblarg\leq@section@nostar}%
    }
    \def\leq@section@star#1{%
      \global\LeqStarSectiontrue
      \gdef\leq@current@section{#1}%
      \gdef\leq@current@subsection{}%
      \leq@original@section*{#1}%
      \global\LeqStarSectionfalse
    }
    \def\leq@section@nostar[#1]#2{%
      \global\LeqStarSectionfalse
      \gdef\leq@current@section{#2}%
      \gdef\leq@current@subsection{}%
      \leq@original@section[#1]{#2}%
    }
    % Redefinir \subsection
    \renewcommand{\subsection}{%
      \@ifstar{\leq@subsection@star}{\@dblarg\leq@subsection@nostar}%
    }
    \def\leq@subsection@star#1{%
      \global\LeqStarSectiontrue
      \gdef\leq@current@subsection{#1}%
      \leq@original@subsection*{#1}%
      \global\LeqStarSectionfalse
    }
    \def\leq@subsection@nostar[#1]#2{%
      \global\LeqStarSectionfalse
      \gdef\leq@current@subsection{#2}%
      \leq@original@subsection[#1]{#2}%
    }

    % ----- Restauración de valores Globales de estilo ----------%
    % Flags
    \newif\ifInFootnote
    \InFootnotefalse
    \pretocmd{\@footnotetext}{\InFootnotetrue}{}{}
    \apptocmd{\@footnotetext}{\InFootnotefalse}{}{}
    % Profundidad de listas (soporta anidadas)
    \newcount\MyListDepth \MyListDepth=0
    \AtBeginEnvironment{la}{\advance\MyListDepth by 1}
    \AfterEndEnvironment{la}{\advance\MyListDepth by -1}
    \AtBeginEnvironment{lnum}{\advance\MyListDepth by 1}
    \AfterEndEnvironment{lnum}{\advance\MyListDepth by -1}
    \AtBeginEnvironment{itemize}{\advance\MyListDepth by 1}
    \AfterEndEnvironment{itemize}{\advance\MyListDepth by -1}
    \AtBeginEnvironment{enumerate}{\advance\MyListDepth by 1}
    \AfterEndEnvironment{enumerate}{\advance\MyListDepth by -1}
    % Restauración diferida: hazlo al INICIO del próximo párrafo real
    \newif\if@pendingRestore
    \@pendingRestorefalse
    \newcommand{\RequestRestore}{\global\@pendingRestoretrue}
    \everypar{%
      \if@pendingRestore
        \global\@pendingRestorefalse
        \ifInFootnote\else
          \ifnum\MyListDepth=0
            \setlength{\parskip}{\GlobalParskip}%
            \setstretch{\GlobalBaseLineStretch}%
          \fi
        \fi
      \fi
    }
    % Al final de bloques:
    \AfterEndEnvironment{equation}{\RequestRestore}
    \AfterEndEnvironment{equation*}{\RequestRestore}
    \AfterEndEnvironment{la}{\RequestRestore}
    \AfterEndEnvironment{lnum}{\RequestRestore}
    % Al final de macros:
    \apptocmd{\eqblock}{\RequestRestore}{}{}
    \apptocmd{\imagenfit}{\RequestRestore}{}{}
    
\makeatother

% ===================== ToC propio con 4 niveles (único bloque) =====================
\makeatletter

% --- Escritores: añadimos una línea al .stc en cada cabecera numerada ---
% Guardamos las versiones actuales para llamar al comportamiento normal
\let\MyTOC@orig@section\section
\let\MyTOC@orig@subsection\subsection
\let\MyTOC@orig@subsubsection\subsubsection
\let\MyTOC@orig@paragraph\paragraph

% SECTION
\renewcommand{\section}{%
  \@ifstar{\MyTOC@section@star}{\@dblarg\MyTOC@section@nostar}%
}
\def\MyTOC@section@star#1{%
  \MyTOC@orig@section*{#1}% sin entrada al .stc
}
\def\MyTOC@section@nostar[#1]#2{%
  \MyTOC@orig@section[{#1}]{#2}% comportamiento normal (escribe al .toc)
  % Escribimos también nuestra línea al .stc (texto con \numberline)
  \addcontentsline{stc}{section}{\protect\numberline{\thesection}#2}%
}

% SUBSECTION
\renewcommand{\subsection}{%
  \@ifstar{\MyTOC@subsection@star}{\@dblarg\MyTOC@subsection@nostar}%
}
\def\MyTOC@subsection@star#1{%
  \MyTOC@orig@subsection*{#1}%
}
\def\MyTOC@subsection@nostar[#1]#2{%
  \MyTOC@orig@subsection[{#1}]{#2}%
  \addcontentsline{stc}{subsection}{\protect\numberline{\thesubsection}#2}%
}

% SUBSUBSECTION
\renewcommand{\subsubsection}{%
  \@ifstar{\MyTOC@subsubsection@star}{\@dblarg\MyTOC@subsubsection@nostar}%
}
\def\MyTOC@subsubsection@star#1{%
  \MyTOC@orig@subsubsection*{#1}%
}
\def\MyTOC@subsubsection@nostar[#1]#2{%
  \MyTOC@orig@subsubsection[{#1}]{#2}%
  \addcontentsline{stc}{subsubsection}{\protect\numberline{\thesubsubsection}#2}%
}

% PARAGRAPH
\renewcommand{\paragraph}{%
  \@ifstar{\MyTOC@paragraph@star}{\@dblarg\MyTOC@paragraph@nostar}%
}
\def\MyTOC@paragraph@star#1{%
  \MyTOC@orig@paragraph*{#1}%
}
\def\MyTOC@paragraph@nostar[#1]#2{%
  \MyTOC@orig@paragraph[{#1}]{#2}%
  % Si no numeras los \paragraph, cambia \theparagraph por texto plano sin \numberline
  \addcontentsline{stc}{paragraph}{\protect\numberline{\theparagraph}#2}%
}

% --- Impresor del índice propio (lee el .stc) ---
\newcommand{\MyTOC}{%
  \par\bigskip
  {\centering\bfseries\Large Índice de contenido\par}%
  \medskip
  \begingroup
    \parindent=0pt\parskip=2pt
    \@starttoc{stc}%
  \endgroup
}

\makeatother
% ===================== fin ToC propio =====================


%%%%%%%%%%%%%%


% --- Datos del tema ---
\title{Unión Europea (V): Cohesión Económico y Social en la UE\\
\large Política regional e instrumentos presupuestarios y financieros. Política social y de empleo. El proceso de convergencia real en la Unión Europea}
\author{Marcelo Ortega}
\date{Marzo 2026}
\newcommand{\TemaCode}{3B40}

\oldtitle{
    B.41. La Unión Europea y la Cohesión Económica y Social: política regional y reforma de los fondos estructurales. Política social y de empleo. Implicaciones sobre el proceso de convergencia real en la Unión Europea
}
\changerationale{
    Ajuste de redacción
}

\begin{document}


\maketitle
\changebox

\newpage

% --- Página de resumen y modificaciones ---
\puntosclave{ %% Escribir puntos clave Apuntes CECO
     
    }
\vspace{1cm}

\modificaciones{
    \textbf{NOTA (04/10/2026):} Revisar: https://www.bde.es/wbe/es/publicaciones/analisis-economico-investigacion/documentos-ocasionales/europa-en-transicion-energetica-respuestas-y-desafios-tras-dos-crisis-consecutivas.html
}

\newpage
\MyTOC
\newpage

\mylistofequations
\clearpage

\MyListOfFigures
\clearpage


\pagenumbering{arabic}
\setcounter{page}{1}


\section{Introducción}



\subsection{Contextualización}


\subsubsection{Relevancia}




\subsubsection{Problemática}



\subsection{Estructura de la exposición}









\clearpage
\section{Unión Europea: Política Regional y de Cohesión}
\puntosclave{
    Fortalecer su cohesión económica, social y territorial es uno de los objetivos principales de la Unión. Desde 1988, ha registrado un enorme aumento de su presupuesto y se ha convertido, junto con la PAC, en una de más significativas en términos cuantitativos. 
    
    La Unión apoya la consecución de estos objetivos mediante la utilización de los Fondos Estructurales y de Inversión Europeos.
    
    El Acuerdo de Asociación es el documento nacional, de carácter estratégico, elaborado por cada Estado miembro, que expone la estrategia y prioridades de inversión de los Fondos de la Política de Cohesión.  
    
    Ha habido una convergencia real entre las regiones de la Unión Europea en los últimos años, con una reducción del $34\%$ en la brecha económica entre las regiones más ricas y las más pobres desde 2000, y un aumento más rápido del PIB per cápita en las regiones más pobres. 
    
    Sin embargo, esta convergencia no ha sido uniforme en todas las regiones de la UE, y las regiones del sur y del este de Europa todavía tienen un nivel de desarrollo inferior al de las regiones del norte y del oeste. 
    
    La cooperación y el intercambio de buenas prácticas entre las regiones han facilitado una mejor utilización de los recursos y un mayor aprendizaje entre ellas. 
}

A fin de promover un desarrollo armonioso en todo su territorio, la Unión fortalece su cohesión económica, social y territorial a través de la Política Regional y de Cohesión (histórica Política de Cohesión territorial). En concreto, busca reducir las disparidades entre los niveles de desarrollo de sus distintas regiones. Se presta especial atención a las zonas rurales, a las zonas afectadas por una transición industrial y a las regiones que padecen desventajas naturales o demográficas graves y permanentes como, por ejemplo, las regiones más septentrionales con escasa densidad de población y las regiones insulares, transfronterizas y montañosas.


Las políticas de cohesión territorial de la Unión Europea tienen una base jurídica que se encuentra principalmente en su normativa originaria, de la que se desprende un corpus normativa derivado a través del cual se materializa, regulando particularmente los fondos destinados a este fin.\footnote{%
    A nivel reglamentario, encontramos el Reglamento (UE) 1303/2013, que establece disposiciones comunes sobre los fondos estructurales y de inversión europeos, lo que regula la asignación y la gestión de los fondos destinados a la cohesión territorial. También, el vigente Reglamento (UE) 2021/1060 del Parlamento Europeo y del Consejo sobre disposiciones comunes para los fondos de la política de cohesión para el período 2021-2027. Este reglamento establece los principios generales y las normas para la gestión de los fondos europeos, como el FEDER, el Fondo Social Europeo, el Fondo de Cohesión y el Fondo Europeo Agrícola de Desarrollo Rural, todos con un enfoque claro en la cohesión económica, social y territorial.
} En concreto, el art. 3 del TUE establece que la Unión debe promover una economía social de mercado altamente competitiva, que tenga como objetivo el pleno empleo y el progreso social, y una cohesión económica, social y territorial equilibrada. Por su parte, los artículos 174 a 178 del TFUE resultan fundamentales para las políticas de cohesión territorial ya que establecen los objetivos de éstas: \textbf{reducir las disparidades entre las distintas regiones de la Unión, apoyando el desarrollo y la convergencia económica}. En particular, el artículo 174 hace referencia a la necesidad de fortalecer la cohesión económica, social y territorial, centrándose especialmente en las regiones menos desarrolladas.

No obstante, la legislación en este ámbito está en constante evolución, reflejando el compromiso de la Unión Europeo con una cohesión económica, social y territorial más equilibrada.

\textbf{¿Qué queremos decir?} A lo largo de esta primera sección, veremos....

De acuerdo con el Tratado de Funcionamiento de la Unión Europea (TFUE), la cohesión económica, social y territorial es una política compartida, lo que implica que tanto la UE como los Estados miembros pueden legislar y adoptar actos jurídicamente vinculantes en dicho ámbito, y que estos últimos ejercerán su competencia sólo en la medida en que la Unión no haya ejercido la suya. Además, establece que la Unión Europea, con el fin de promover un desarrollo annonioso del conjunto de la Unión, reforzará su cohesión económica, social y territorial, reduciendÓ las diferencias entre los niveles de desarrollo de las diversas regiones.

\subsection{Desarrollo histórico: ¿Qué persigue esta política comunitaria?}
La evolución histórica de las políticas de cohesión territorial se remonta a la mismísima concepción de la Comunidad Económica Europea, en 1957.

El objetivo de la Política Regional Europea (hoy, Política de Cohesión Económica y Social) ha sido, desde sus inicios, reducir las disparidades en el nivel de renta y riqueza entre las distintas regiones (y Estados miembros) que la integran. Para ello, a lo largo de los sucesivos períodos de programación, ha movilizado un importante volumen de recursos a través de sus instrumentos: los Fondos Estructurales y el Fondo de Cohesión, con el fin de promover el desarrollo de las regiones más desfavorecidas y avanzar en la corrección de sus deficiencias en cuanto a la dotación de determinados factores productivos (capital físico y capital humano) y mejorar el nivel de vida de sus habitantes a través del aumento de su renta per cápita.

La Unión Europea (UE) ha recorrido un largo camino de más de cincuenta años en los que fue avanzando en el proceso de integración. Las sucesivas adhesiones también supusieron un esfuerzo extraordinario de acoplamiento. Para ajustarse a las exigencias que iban surgiendo  se vio en la necesidad de realizar un proceso importante de reformas. Entre ellas, definir las  medidas necesarias conducentes a reducir las disparidades regionales en todo su ámbito y con un extraordinario esfuerzo de solidaridad ha conseguido ese anhelado sueño de permitir que la Cohesión Económica y Social se haya convertido en uno sus objetivos más importantes.

El avance en el proceso de integración europea y las sucesivas ampliaciones a que se ha visto sometida han permitido una nueva configuración de las disparidades regionales. Aunque  se ha avanzado en su disminución, las zonas menos prósperas de la UE han tenido que hacer frente al aumento de las desigualdades en su interior. Su debilidad procede fundamentalmente de su limitada dotación en capital humano cualificado, de su deficiente dotación de infraestructuras de telecomunicaciones y de su insuficiente inversión en investigación y desarrollo.

Las inversiones realizadas en torno al capital físico se han adjudicado con el fin de modernizar la estructura productiva de las empresas de las regiones menos avanzadas y también para dotarlas de mejores infraestructuras viales (construcción de autopistas, puentes, aeropuertos, las redes ferroviarias, etc.). El capital humano, es decir, todos ciudadanos del ámbito de la UE se han beneficiado de los recursos destinados a través del Fondo Social Europeo (FSE) para mejorar su formación, educación y experiencia, facilitándoles así el acceso al mercado de trabajo.

Los instrumentos utilizados para la consecución de sus objetivos son los Fondos Estructurales y el Fondo de Cohesión. Los Fondos Estructurales están integrados por el Fondo Europeo de Desarrollo Regional (FEDER), El Fondo Europeo de Orientación y Garantía Agrícola en su sección Orientación (FEOGA-O) y el Fondo Social Europeo (FSE). Cada uno de ellos se fue ajustando a los objetivos planteados en cada uno de los periodos de programación.

Otro de los elementos que influyeron, de manera importante, en el proceso de transformación de la UE fueron los efectos profundos de las actuales tendencias sociales y económicas que surgieron por la globalización de la economía y que condujeron a la transformación substancial de la economía europea hacia actividades basadas en la competitividad, productividad y la innovación.



\subsubsection{Orígenes: Comunidad Económica Europea (1957)}
La Política Regional en la Unión Europea actual Política de Cohesión Económica y Social no se concebía como una prioridad en el Tratado de Roma (Tratado CEE, enero de 1958).\footnote{%
    Sin embargo¸ en el preámbulo del Tratado Constitutivo (marzo de 1957) se incluye como uno de los objetivos básicos de los países integrantes \textit{reforzar la unidad de sus economías y asegurar el desarrollo armonioso reduciendo las diferencias entre las diversas regiones y el retraso de las menos favorecidas}. Igualmente, y de forma muy general, el artículo 2 del Tratado incluía como objetivo \textit{promover un desarrollo armonioso de las actividades económicas en el conjunto de la Comunidad, una expansión continua y equilibrada, una estabilidad creciente, una elevación acelerada del nivel de vida y relaciones más estrechas entre los Estados que la integran}. En el Tratado Constitutivo se aborda la Política Regional de forma muy incipiente, pero deja abierta la posibilidad de que los seis países fundadores (Alemania, Francia, Italia, Bélgica, Holanda y Luxemburgo) pudiesen reducir sus desequilibrios a través de intervenciones de carácter nacional.
}

Cuando el Tratado entró en vigor, los instrumentos financieros disponibles para acciones estructurales se limitaban al Fondo Social Europeo (FSE, 1958) y el Banco Europeo de Inversiones (BEI, 1958). La Política Social, adoptada entonces, facilitaba la inserción de la mano de obra entre las diferentes regiones (facilitando la movilidad geográfica y ocupacional). Por otro lado, al instrumento financiero o de crédito se le asignó la función de contribuir al desarrollo equilibrado y estable del Mercado Común a través del fomento del desarrollo regional, la modernización del sector industrial y el manejo racional de la energía. Estas actuaciones se dirigían a entidades públicas y privadas por medio de la concesión de préstamos para la realización de obras de infraestructuras y equipamiento productivo. Paralelamente, la acción de la PAC, instrumentada a través del FEOGA, se considereba también como una política de carácter cohesionador y de convergencia.

En definitiva, aunque en el Tratado de Roma se observan claras intenciones sobre la implementación de la Política Regional, no se adopta de forma oficial. Pero esto no debe verse como una procrastinación o indiferencia injustificada, debemos tener en cuenta que la Comunidad Económica Europea nacía en el seno de unos países con un nivel de vida y de infraestructuras muy  similares. Además, tenían un nivel de desarrollo muy homogéneo, con excepción de Italia (especialmente en el sur del país) con un PIB per cápita muy inferior a la media de la CEE y un elevado nivel de desempleo. Su integración no supuso entonces ningún problema debido a la coyuntura económica favorable y a la intervención oportuna del FSE y los recursos financieros del BEI.\footnote{%
    Asimismo, el Tratado de Roma se basó en los principios neoclásicos según los cuales se confiaba en la capacidad del mercado como un factor importante para la corrección de los desequilibrios regionales.
}


\subsubsection{Interludio (1960-1970): La construcción del sistema}
La década de los sesenta estuvo caracterizada por un gran auge económico que permitió a los Estados miembros emprender una senda de notable crecimiento económico, pleno empleo y completa movilidad de los factores productivos. Este escenario favorable consolidan la convicción de que se generaría en todo su ámbito un importante desarrollo interregional, por lo que los desequilibrios regionales que se fueron gestando en su interior carecieron de importancia. No obstante, los diferentes gobiernos adoptaron una Política Regional individualizada, nacional y sin coordinación comunitaria. 

La situación generó que los desequilibrios regionales existentes se agudizaran y, en otras, persistiesen. Además, se constató que la adopción de la Política Agrícola Comunitaria (PAC) había acentuado aún más las diferencias existentes entre las regiones prósperas y las menos avanzadas. Todo esto produjo que, en 1965, la Comisión Europea reconociese la necesidad de coordinar las intervenciones para solucionar los desequilibrios regionales.\footnote{%
    De esta primera reflexión fue fruto la primera comunicación sobre Política Regional a partir de las conclusiones de tres grupos de expertos, dando lugar en 1968 a la creación de la Dirección General de Política Regional. Poco después se produjo la primera propuesta del Consejo y se instrumentaron una serie de acciones canalizadas hacia el ámbito regional. Su objetivo estaba dirigido hacia las regiones menos avanzadas. Se propone entonces la creación de un Comité Permanente, con la función de ejecutar el plan en mención y poner en funcionamiento un Fondo de Apoyo que para tal efecto se asignó en el BEI.
} Las sucesivas propuestas, tanto de la Comisión como del Consejo, fueron aprobadas por el Comité Económico y Social y el Parlamento Europeo, dando inicio a la Política Regional de los seis países fundadores por vía administrativa.

A lo largo de la siguiente década, la idea de la ayuda estructural para las regiones menos favorecidas empieza a tomar forma.\footnote{%
    De hecho, desde otro ámbito de actuación pero también especialmente importante dada su dotación, en 1971 se autorizan incentivos al desarrollo regional en la Política Agrícola Comunitaria (PAC), ya quese coordinan las ayudas financieras con la intervención del Fondo Europeo de Orientación y Garantía Agrícola en su sección Orientación (FEOGA-O) en actuaciones de desarrollo regional. Esto fue motivado, fundamentalmente,  gracias a varias Resoluciones del Consejo Europeo.
} Tan pronto como 1972, en la Conferencia de jefes de Estado (realizada en París), se decide que para reforzar la Comunidad era necesario definir una verdadera Política Regional. 

Con ese objetivo se realiza el primer gran informe sobre la materia: Informe THOMPSON (Comisión, 1973). En éste, se hace un diagnóstico de la situación de las regiones que integraban la entonces CEE y se proponen varias directrices para la corrección de los desequilibrios regionales observados.\footnote{%
    El informe destaca la importancia de la Política Regional comunitaria para reducir las distorsiones estructurales, pues se destaca que eran el \textbf{mayor obstáculo} para la consecución, en el futuro, \textbf{de la Unión Monetaria}. También se insiste en la importancia de que la intervención comunitaria no debía sustituir en el futuro a las políticas nacionales sino que debían ser \textbf{complementarias}. El Informe THOMSON también se propone la creación del Fondo Europeo de Desarrollo Regional (FEDER) y el Comité de Política Regional, como un órgano de carácter institucional con la misión de coordinar la política regional de los Estados Miembros.
} Además, en 1973 ingresan el Reino Unido, Dinamarca e Irlanda; y, con esta última, se acentuaron las diferencias regionales: Irlanda aportaba a la Comunidad la renta per cápita más baja de la media comunitaria. Igual sucedía con Italia que venía arrastrando desajustes estructurales graves (desequilibrios regionales que no habían cedido después de un largo periodo de integración en la CEE). Puestas en conjunto, esta coyuntura incidió para que Italia e Irlanda presionaran sobre la necesidad de que les fuesen transferidos fondos para atenuar sus disparidades y alcanzar cierta homogeneidad con sus socios comunitarios.

Las circunstancias planteadas permitieron el nacimiento de la Política Regional de la CEE (1975) y, en este mismo año y después de muchas controversias, también del Fondo Europeo de Desarrollo Regional (FEDER, 1975).\footnote{%
    No obstante, su dotación inicial fue irrisoria, pues equivalía al $4,8\%$ del total del Presupuesto Comunitario, lo que era manifiestamente insuficiente para hacer frente a los problemas estructurales existentes. Con esto vemos como, obviamente, la solución de los desequilibrios regionales no era una prioridad en la CEE en aquella época, ya que tampoco se tuvo en cuenta que varios países (Irlanda e Italia, especialmente) estaban inmersos en una crisis económica y energética de la que tardarían varios años en salir. A pesar de esto, y dadas la circunstancias, se decide poner en funcionamiento el Comité de Política Regional (creado por la Decisión del Consejo 185/75/CEE de 18 de marzo), como Órgano Consultivo del Consejo y la Comisión, con carácter institucional, con la misión de definir los problemas regionales existentes en la CEE y establecer, a partir de los mismos, los mecanismos necesarios para la corrección de los desequilibrios regionales existentes. De este nuevo órgano emanan las directrices de funcionamiento del FEDER que se exponen seguidamente ($75\%$).
} El FEDER empieza a funcionar como un instrumento destinado a financiar proyectos de inversión, especialmente en regiones deprimidas con una renta inferior al $75\%$ de la media de los nueves estados miembros e interviene, especialmente, para \textbf{paliar los problemas de las zonas industriales en declive industrial}. En esta primera etapa, la intervención del Fondo cumple un papel exclusivamente redistributivo, debido a la inexistencia de objetivos comunes y la ausencia de criterios compartidos en la selección de los proyectos o en la elección de las regiones. A pesar de ello, el FEDER fue creciendo gracias a las luchas presupuestarias que se suscitaron entre el Consejo y el Parlamento Europeo. Este último defendía los Gastos No Obligatorios del Presupuesto Comunitario en los que estaban incluidos los Fondos Estructurales (el FEDER, el FSE y el FEOGA).

A finales de la década de los setenta, a pesar de las medidas adoptadas hasta entonces en materia de desarrollo regional, las diferencias regionales en el seno de la Comunidad no habían disminuido. Al contario, las diferencias en términos de renta per cápita seguía aumentando su disparidad.\footnote{%
    En vista de lo anterior, la Comisión elaboró las nuevas orientaciones en materia de Política Regional, que presentó al Consejo en comunicación de 3 de junio de 1977. Asimismo, remitió al Consejo varias propuestas, entre ellas, la modificación del Reglamento con el que se creaba el FEDER, la elaboración de un informe bianual sobre la evolución socioeconómica de las regiones (informe que debería ser aprobado por el Consejo a propuesta de la Comisión) y, por último, el procedimiento mediante el cual se asignaba la dotación del FEDER.
} Como consecuencia, a partir de diversas propuestas planteadas por la Comisión, se modificó en 1979 el FEDER (Reglamento 214/79/CEE), en vigor a partir de 1980. 

Con el nuevo reglamento se consolidó la situación del Fondo y que su dotación anual vendría fijada en el Presupuesto General de la Comunidad. A partir de entonces el FEDER cuenta con un apoyo legal a través del cual se sentaron las bases para \textbf{acciones concretas en áreas y sectores diferentes}. Con esta reforma de los Reglamentos del FEDER, se fue consolidando y perfilando con algunas deficiencias lo que sería a partir de entonces una verdadera política de desarrollo regional la de la CEE.

\textcolor{red}{se dieron los nuevos pasos hacia un avance importante de la Política Regional comunitaria; se le concedieron al Fondo facultades para orientar sus recursos en función del conjunto de países beneficiarios y no de forma individualizada como se venía haciendo. En definitiva,}

\subsubsection{Impulso y consolidación (1980-2010): Las importante reformas y adopción de la Política Regional}
La CEE vivió desde comienzos de los años ochenta en clima de reformas internas con el fin de superar el estancamiento que venía arrastrando desde la década de los setenta. Como consecuencia de la situación se produce una nueva orientación en la Política de Desarrollo Regional, consignada en la proposición del Reglamento, presentada por la Comisión al Consejo el 29 de octubre de 1981 (J.O. número C 336 de 23 de diciembre de 1981), que modifica el Reglamento 724/75 (del que hemos venido haciendo referencia). 

Las nuevas orientaciones de la Comisión partían de la base de que para aumentar el impacto de las actuaciones del FEDER era necesario llevar a cabo algunas modificaciones motivadas por la insuficiencia de recursos para hacer frente a la corrección de los desequilibrios regionales existentes. Se propuso concentrar las intervenciones en regiones afectadas seriamente por un subdesarrollo estructural, buscando que el FEDER interviniese en las regiones más desfavorecidas \textbf{considerando la Comunidad en su conjunto y no de forma individual} como hasta entonces se había venido haciendo. Aunque con esta medida muy pocos Estados miembros resultarían beneficiados (algunas zonas de Grecia, Italia, Reino Unido e Irlanda), fue aprobada por el Consejo con algunas modificaciones. Es así como el Reglamento 724/75 pasó a ser reemplazado por el 1787/84, en vigor a partir de enero de 1985.\footnote{%
    El nuevo reglamento del FEDER (1787/84) supuso un gran avance en la adopción de la Política Regional. Le concede mayor autonomía en sus actuaciones con el fin de reducir los desequilibrios regionales de forma más efectiva. El Reglamento constaba de cuatro títulos: coordinación de las políticas regionales de los Estados miembros, disposiciones generales relativas del Fondo, disposiciones particulares sobre las intervenciones y disposiciones sobre compromisos y pagos. Hace hincapié en el impacto de la intervención sobre el desarrollo regional, buscando siempre redistribuir de la forma más equilibrada posible su actividad económica. Se incluyeron zonas geográficamente determinadas y afectadas con problemas especialmente graves (renta per cápita muy baja, deficiencias en infraestructuras, etc.).
}

Las aportaciones del FEDER se destinaron a financiar tres objetivos: \textbf{desarrollo industrial, nuevos servicios e infraestructuras de comunicaciones}.\footnote{%
    También se concedió especial importancia a las pequeñas y medianas empresas, los artesanos y al sector turístico. Todas estas actuaciones tenían como objetivo incrementar su actividad productiva a través de la adopción de nuevas y mejores técnicas de producción, favorecidas por la incorporación de innovaciones tecnológicas con el fin de mejorar la competitividad.
} Además, la Comisión sería el órgano facultado para aprobar la concesión de las ayudas y supervisar su ejecución presupuestaria, que se adjudicarían de forma coherente con los programas de desarrollo regional y se debían ajustar a los objetivos y prioridades planteadas por Comisión.

No obstante, las sucesivas adhesiones aumentaron las desigualdades espaciales haciéndola cada vez más heterogénea (primero Irlanda, después Grecia y más adelante España y Portugal) y en un contexto de profunda crisis económica. En este sentido, las reformas del FEDER hicieron su operatividad muy difícil creando expectativas inalcanzables. Muchas de ellas se sucedieron de forma muy rápida y en la mayoría de los casos se adoptaron decisiones políticas con el único fin de cumplir con compromisos de diversa índole y no en función de las necesidades regionales. De igual forma, los recursos asignados, aunque crecieron de forma notable, no evolucionaron de acuerdo con la magnitud y variedad de los problemas existentes.\footnote{%
    Esta situación de descoordinación tal, provcó que se adoptase de manera piloto el Enfoque Integrado en la aplicación de los Fondos Comunitarios con finalidad estructural. Las acciones definidas fueron: los Programas Integrados Mediterráneos (PIM), los Programas Integrados de Desarrollo (PID) y las Operaciones Integradas de desarrollo (OID). Con la adopción de estas medidas se buscaba coordinar, en sus actuaciones conjuntas, la intervención del FEDER, el FSE y el FEOGA-O hacia una zona afectada por una situación socioeconómica grave. En definitiva, esta iniciativa no permitió alcanzar el objetivo deseado y en lugar de simplificar las actuaciones lo que hizo fue hacerlas más complejas. Se produjo una gran confusión en la adopción de las distintas  medias y su actuación se asumió como una simple experiencia piloto.
}

\paragraph{Acta Única Europea (1985): Una reforma multidimensional}
Tras estos intentos, y con un rejuvenecido impulso comunitario, la transformación definitiva de la Política de Cohesión se inicia en 1985. Como sabemos, en junio de 1985 se convoca una Conferencia Intergubernamental con el fin de proceder a una serie de modificaciones institucionales en el ámbito comunitario y la extensión de actividades con miras a la creación del Mercado Único y la futura Unión Europea, de la cual nace la famosa Acta Única Europea, AUE, (Milán, junio de 1985). El AUE es un avance importante en el proceso de integración que pone de manifiesto la necesidad de iniciar la consecución del Mercado Interior (MI) o el Mercado Único Europeo (MUE): un gran espacio sin fronteras que diera lugar a la libre circulación de bienes, capitales, servicios y personas [\authorfont{Ver Tema 3.B.39}]. Además, en el Libro Blanco sobre el Mercado Interior se identifican los obstáculos que impedían su consecución: físicos, técnicos, fiscales y, relevante para nuestro objeto, las dificultades que supondría para las economías más débiles su funcionamiento.\footnote{%
    Se argumentaba que las inversiones se trasladarían a las regiones más desarrolladas en detrimento de las menos avanzadas y que inevitablemente esta situación conduciría al incremento de las disparidades regionales, poniendo en peligro el objetivo de convergencia económica. En este contexto, la Comisión insiste en la importancia de alcanzar la Cohesión Económica y Social como sostén de la Política Regional. Además, insiste en los efectos negativos que se producirían tras el establecimiento de un gran mercado si su puesta en práctica no se acompañaba de las medidas necesarias para desarrollar las regiones con retraso estructural. Para ello, se propone mejorar las condiciones de empleo, estimular la innovación social, reconvertir las zonas industriales en declive, etc. De igual manera la Comisión determina que los instrumentos financieros existentes no contaban entonces con la eficacia correctora necesaria para la superación de los desequilibrios que se venían arrastrando y los que se sumarían después de la cuarta ampliación. Esta iniciativa fue defendida por Grecia e Irlanda y apoyada por España y Portugal (candidatos entonces).
} 
    
Con esto en mente, el AUE (art. 130D) recoge la propuesta de \textit{introducir en la estructura y en las normas de funcionamiento de los Fondos con finalidad estructural las modificaciones que fuesen necesarias para precisar y racionalizar sus funciones, reforzar su eficacia y coordinar sus intervenciones entre sí con las de los instrumentos financieros existentes}. Por tanto, establece la Política Regional como uno de los \textbf{objetivos prioritarios de la CEE}. Se adopta entonces un bloque de medidas, entre ellas una nueva reforma de los Fondos Estructurales, a los que se les concede a partir de entonces la misión de servir como instrumentos de desarrollo y se otorga a la Comunidad la responsabilidad de reducir las disparidades regionales existentes en su ámbito.\footnote{%
    A partir de esta propuesta se origina el Reglamento 2052/88 del Consejo de junio de 1988, también denominado Reglamento Marco, por el que se regulaba, junto con el Reglamento de coordinación CEE número 4253, la financiación de naturaleza estructural durante el periodo 1988-1993, que estarían integrados por el FEDER, el FSE y el FEOGA-O. El FEDER y participarían en el desarrollo y ajuste estructural de las regiones menos avanzadas (con una renta per cápita inferior al 75 por 100 de la media de la CEE=12) y en la reconversión de las regiones industriales en declive. Los demás instrumentos (FSE, FEOGA-O y el BEI) apoyarían al FEDER en sus intervenciones.
}
    
La reforma concede especial atención a la concentración de las intervenciones financieras en un limitado número de objetivos, el establecimiento de un nuevo enfoque de aplicación y gestión de los mismos y un incremento importante de las aportaciones. Además, incorpora los cuatro grandes principios generales que regulan el funcionamiento de los Fondos Estructurales: (i) concentración, de las intervenciones destinadas a financiar proyectos en las regiones menos avanzadas; (ii) cooperación, entre la Comisión y las autoridades de los Estados miembros (nacionales regionales y locales); (iii) programación, programación plurianual de las intervenciones a través de la elaboración de programas que se ejecutarán durante varios años; y, (iv) adicionalidad, las aportaciones comunitarias no deben sustituir a los gastos nacionales, deben ser complementarios (cofinanciación). Además, para la gestión de los Fondos \textbf{se adopta entonces la nomenclatura común de las Unidades Territoriales Estadísticas (NUTS)}, incorporadas por Eurostat con el fin de ofrecer una división uniforme de unidades territoriales para la elaboración de las estadísticas regionales de la Unión Europea. 

Con los objetivos plasmados en el Acta Única Europea, las acciones estructurales alcanzaron una nueva dimensión. Fue necesario recorrer un largo camino en el que la Política Regional se fuese haciendo un sitio importante en el conjunto de políticas de la Unión Europea. Sin embargo, las expectativas del Mercado Único incidieron de forma positiva y determinante para que la Cohesión Económica y Social se convirtiese en uno de sus objetivos más importantes. 

\begin{specialblock}{Reformas: Particularidades y nueva estructura}
    En el Reglamento recién aprobado recogía un incremento importante en la cuantía de recursos asignados para financiar la Política Regional (para los años 1987-1993 se duplicaron en términos reales). La clasificación de las regiones y la intervención de los Fondos se estructuran de la siguiente manera:\footnote{%
        Existe más objetivos, pero son menos relevantes: el objetivo número 5a con la misión de facilitar la adaptación de las estructuras agrarias, se financiaba con cargo al FEOGA-O; y el objetivo número 5b permitía la modernización y comercialización de su producción e intervenían el FEOGA-O, el FEDER y el FSE.
    }
    \begin{lnum}
        \item \textbf{Prioritarias (I): Desarrollo}. Las regiones incluidas en el objetivo número 1, con una renta per cápita inferior al 75 por 100 de la media de la CEE=12, serían las beneficiadas con los recursos del FEDER;
        \item \textbf{Prioritarias (II): Coyuntural}. Las regiones incluidas en el objetivo número 2, deberían presentar un porcentaje de paro industrial superior al de la media comunitaria con disminución del empleo desde 1975 (regiones industriales en declive). Se financiarían con recursos del FEDER y el FSE;
        \item \textbf{Complementarias: Comunitarias}. En los objetivos números 3 y 4 no se establece delimitación alguna, es decir, podrían beneficiarse todas las regiones de la CEE. El objetivo 3 con recursos para favorecer a los parados de larga duración y el 4 para financiar la inserción profesional de los jóvenes. Los dos objetivos se financian con cargo al FSE.
    \end{lnum}
    
    La Comisión insiste en la necesidad de reforzar la colaboración con los Estados miembros beneficiados de las ayudas y en el control y gestión de los recursos. La complementariedad adquiere, por tanto, gran importancia ya que la intervención comunitaria se realizaría a partir de entonces teniendo en cuenta las necesidades propuestas por los Estados miembros. También se crea un nuevo instrumento de solidaridad: las Iniciativas Comunitarias, con el fin de solucionar situaciones difíciles relacionadas con otras políticas comunitarias y a través de ellas, facilitar la consecución de sus objetivos.

    La novedad de esta reforma estriba en que, por primera vez, se incluyen en un marco legal todos los medios operativos básicos con los que contaba la Política Regional para alcanzar sus diferentes objetivos. Fue un paso importante para su estructura organizativa, al asignar los diferentes Fondos Estructurales a los objetivos establecidos en el Acta Única Europea. Con lo cual la reforma permitió potenciarlos y reorientarlos de tal manera que fuera posible el objetivo de Cohesión Económica y Social en todo su ámbito y facilitar el Mercado Único.
\end{specialblock}

\paragraph{Tratado de Maastricht (1992): Consolidación}
Poco después, se vuelve adar un nuevo impulso a la Política de Cohesión a través del Tratado de Maastricht, donde se modifican los fundamentos de la integración: dar paso a una Unión Económica, definida como un Mercado Único, políticas comunes y política económica común (especialmente la monetaria). 

Gracias a esta reforma y clara definición de objetivos, se refuerza la Cohesión Económica y Social articulándose como uno de sus objetivos prioritarios (arts. 158 a 162), buscando con ello alcanzar un mayor equilibrio entre los niveles de crecimiento económico en las regiones que integran la Unión Europea.\footnote{%
    Además, también en el Tratado de Maastricht, se introduce el Principio de subsidiariedad que, unido a los cuatro previamente incluidos: concentración, cooperación, programación y adicionalidad, facilitaron el proceso de programación en la concesión de las ayuda dentro de la Política Regional de la UE.
} Para cumplir con los objetivos se incrementa la dotación de recursos, que ascendían a la cuarta parte del Presupuesto Comunitario en DELORS I (en detrimento de la PAC, que a partir de entonces empieza a recibir menos aportaciones). Además, se refuerza mediante la creación de numerosos instrumentos de diversa índole. En el ámbito financiero se crea el Fondo de Cohesión (1993) y el Instrumento Financiero de Orientación Pesquera (IFOP), con el fin de contribuir a la realización de los objetivos de la Política Pesquera Común. En el ámbito político, se crea el Comité de las Regiones, una medida innovadora que facilitó, a partir de entonces, la comunicación entre las instituciones de la UE y los diferentes estamentos nacionales, regionales y locales, posibilitando la consecución de los objetivos previstos en el terreno regional.\footnote{%
    Una vez aprobado el Tratado de Maastricht, la Comisión presentó un nuevo paquete de medidas denominado Paquete Delors II (1992), plasmado los Reglamentos CEE 2081/93 y 2082/93 que incluyen las disposiciones establecidas en el Tratado de Maastricht para el periodo de programación 1994-1999. La novedad del Reglamento es la inclusión de nuevos  mecanismos de evaluación en las intervenciones con el fin de valorar la gestión y el impacto alcanzado con sus intervenciones.
}

En definitiva, podríamos concluir que el periodo que se expande desde 1980 hasta 2010 constituye el periodo de consolidación de la Política Regional como uno de los objetivos prioritarios de la comunidad. Tanto por su influencia en el buen funcionamiento del MUI, como por la redefinición de lo que hasta entonces se consideraba únicamente un proyecto de integración económica, con alguna connotación política, la Política de Cohesión alcanza desde entonces un papel prioritario. De hecho, en linea con la desovietización de las economías del este, la Política de Cohesión se configura como el prinicpal instrumento de convergencia (es decir, de actuación ex ante), que facilitará la incorporación de estas economías a la Unión, así como de nivelación ex post, por las grandes disparidades que presenta la Unión tras la gran ampliación de los primeros años del mileno.\footnote{%
    Ante las perspectivas generadas por el ingreso de nuevo socios y el previsible incremento de los desequilibrios regionales se refuerza la Cohesión Económica y Social. Por ello, se insiste en la importancia de promover el progreso económico y social y un alto nivel de empleo y conseguir un desarrollo equilibrado y sostenible. Tras integrar la UE-$15$ a diez nuevos miembros con un nivel de desarrollo económico significativamente menor, se llegó al acuerdo de adaptar la política regional a la nueva situación, de manera que las regiones de la UE$-15$ siguieran recibiendo recursos y los contribuyentes netos vieran limitada su contribución presupuestaria.
}

\begin{specialblock}{Fondo de Solidaridad (2004): Instrumentos de actuación ex ante}
    El Fondo de Solidaridad es el principal instrumento de la Unión para ayudar a los Estados miembros y a los países candidatos y en vías de adhesión a la Unión a recuperarse de catástrofes naturales graves o emergencias de salud pública. Proporciona ayuda financiera para la reconstrucción tras inundaciones, incendios forestales, terremotos, tormentas o sequías. Desde 2020, también ha cubierto crisis sanitarias graves, como la pandemia de COVID-19. Debido al aumento de la frecuencia y la gravedad de los fenómenos meteorológicos extremos causados por el cambio climático, el papel del Fondo como expresión fundamental de la solidaridad europea ha cobrado cada vez más importancia, mientras que la respuesta rápida a las crisis sigue siendo competencia de la Reserva para Ayudas de Emergencia. El Fondo de Solidaridad fue creado en 2002 para responder a las devastadoras inundaciones que afectaron a Europa central en el verano de ese año.

    El Fondo de Solidaridad ha prestado apoyo en más de 130 ocasiones (en casos de catástrofes naturales e intervenciones en respuesta a emergencias de salud pública). Hasta la fecha, veinticuatro Estados miembros (más el Reino Unido antes de su retirada de la Unión) y cuatro países en vías de adhesión (Albania, Montenegro, Serbia y Turquía) han recibido ayuda del Fondo de Solidaridad y se han pagado en total más de 8 600 millones EUR. Desde 2021, el Fondo de Solidaridad y la Reserva para Ayudas de Emergencia se han financiado como un instrumento, denominado Reserva de Solidaridad y para Ayudas de Emergencia. El presupuesto anual máximo para la Reserva de Solidaridad y para Ayudas de Emergencia es de 1 200 millones EUR (a precios de 2018). En febrero de 2024, se revisó el marco financiero plurianual 2021-2027 y el presupuesto máximo anual para la Reserva de Solidaridad y para Ayudas de Emergencia se incrementó en 1 500 millones EUR con el fin de mejorar la capacidad de la Unión para hacer frente a crisis y situaciones de emergencia. En diciembre de 2024, se aprobó el Reglamento relativo al apoyo regional urgente para la reconstrucción (RESTORE) [Reglamento (UE) 2024/3236], que proporciona apoyo adicional a los Estados miembros que se han visto afectados por una catástrofe natural entre 2024 y 2025 y les permite reprogramar los fondos procedentes del Fondo Europeo de Desarrollo Regional (FEDER), el Fondo de Cohesión y el Fondo Social Europeo Plus (FSE+). Los recursos totales que se pueden reprogramar en el marco del apoyo regional urgente para la reconstrucción están limitados al 10 % de la asignación nacional inicial total para el FEDER y el FSE+. Con ello se ponen a disposición recursos para reparar las infraestructuras dañadas, proporcionar alimentos y asistencia material básica y apoyo social y sanitario.

    La intervención del Fondo de Solidaridad adopta la forma de una subvención que complementa el gasto público del Estado beneficiario y se destina a financiar medidas de emergencia y recuperación esenciales que mitiguen los daños que, en principio, no sean asegurables. Las medidas urgentes que pueden optar a financiación son las siguientes: — restablecimiento inmediato del funcionamiento de las infraestructuras e instalaciones en los sectores de la energía, el agua potable, la eliminación de las aguas residuales, las telecomunicaciones, los transportes, la sanidad y la educación; — puesta a disposición de alojamientos provisionales y financiación de servicios de auxilio destinados a cubrir las necesidades de la población afectada; — aseguramiento inmediato de las infraestructuras de prevención y medidas de protección del patrimonio cultural; — limpieza de las zonas siniestradas, incluidas las zonas naturales; — asistencia rápida, incluida asistencia médica, a la población afectada por una emergencia grave de salud pública y protección de la población frente al riesgo de verse afectada.
\end{specialblock}

\begin{specialblock}{Reformas menores: De Amsterdam (1992) a Niza (2003)}
    En 1999 entra en vigor el Tratado de Ámsterdam, cuyo objetivo es el fortalecimiento de la Europa de los ciudadanos, es decir, la igualdad de oportunidad para los ciudadanos europeos. Cabe destacar que, tan pronto como su segundo artículo señala la Cohesión Social como uno de sus objetivos centrales. En este sentido, consolida los mecanismos establecidos por el Tratado de Maastricht, al enunciar orientaciones sociales prioritarias en el ámbito comunitario, en particular para el empleo.

    Con este Tratado se prepara el camino para iniciar las negociaciones para la futura ampliación hacia la Europa del Este. Para ello, el 26 de marzo de 1999, al final del Consejo Europeo de Berlín, los jefes de Estado y gobierno concluyeron un acuerdo político sobre la \textbf{Agenda 2000}.\footnote{%
        La Agenda 2000 contiene un programa con pautas de actuación cuyo objetivo más importante era dotar a la UE de un nuevo marco financiero para el periodo 2000-2006 de cara a la ampliación. Los aspectos más importantes incluidos en la Agenda 2000 fueron: la reforma agrícola, el aumento de la eficacia en la gestión de los Fondos Estructurales y el refuerzo de la estrategia de preadhesión para los países candidatos a través de varios instrumentos financieros: un instrumento estructural de preadhesión, ISPA, dirigido a mejorar las infraestructuras de transporte y la protección del medio ambiente de los países aspirantes a la adhesión), un instrumento agrícola de preadhesión, SAPARD, que facilitaría la adaptación a largo plazo de las zonas rurales y del sector agrícola de los países candidatos, el programa PHARE con la misión de apoyar financiera y técnicamente a los países de Europa Central y Oriental y el TACIS con la misión de favorecer la transición hacia una economía de mercado y consolidar la democracia y el Estado de Derecho en los Estados socios de Europea Oriental y Asia Central.
    }
    
    Poco después, en la cumbre de Lisboa (23 y 24 marzo de 2000), los Jefes de Gobierno de la UE aprobaron la famosa \textbf{Estrategia de Lisboa}, un nuevo objetivo para la Unión Europea para \textit{convertirla en la economía del conocimiento más competitiva y dinámica del mundo capaz de crecer económicamente de manera sostenible con más y mejores empleos y con mayor Cohesión Social}, antes de 2010. Para la consecución de los objetivos era preciso realizar una serie de reformas dirigidas a potenciar la sociedad de la información, mejorar la educación, realizar reformas económicas, estructurales, dirigidas a estabilizar la moneda y una combinación de políticas macroeconómicas que fuesen favorables al crecimiento económico y a unas finanzas públicas sostenibles (Informe de la Comisión Europea).
    
    En 2003 entra en vigor el Tratado de Niza. Su objetivo se centra en solucionar los \textit{asuntos pendientes} del Tratado de Ámsterdam. En este Tratado se incluyen una serie de reformas institucionales dirigidas a implantar una gestión más eficiente en una futura UE integrada por veintisiete Estados. Entre las principales novedades se incluyen las cooperaciones reforzadas, cuestiones institucionales y la ampliación del uso de mayorías calificadas. El sistema de cooperaciones reforzadas, introducido en el Tratado de Ámsterdam, fue fortalecido en el Tratado de Niza con el fin de facilitar una integración más rápida. La reforma institucional cambió el número de votos asignados en el Consejo Europeo, la asignación de asientos en el Parlamento Europeo y en la Comisión. En la Corte de Justicia, varias de sus características se reforman con el voto unánime del Consejo en lugar de requerirse una Conferencia Intergubernamental. Determinó, asimismo, cambios en la distribución de las competencias entre el Tribunal de Primera Instancia y la Corte de Justicia.
\end{specialblock}





Tras la crisis financiera global de 2008, la política de cohesión adoptó un enfoque más centrado en la eficiencia y la utilización de los fondos para promover la recuperación económica, el empleo y el crecimiento. La Estrategia Europa 2020 marcó el rumbo de la política de cohesión, centrada en la innovación, la inclusión social y la sostenibilidad. Se promovió la cooperación entre regiones para fomentar el crecimiento económico y social.

\subsection{Estructura actual: ¿Cómo se articula la Política de Cohesión?}
Como hemos visto a través de su reccorido histórico, a lo largo de los años la Políticas de Cohesión territorial de la Unión han evolucionado desde un enfoque más básico en la reducción de disparidades económicas hacia un modelo más integral que abarca no solo el desarrollo económico, sino también la innovación, la sostenibilidad, la inclusión social y la competitividad. A continuación presentamos la estructura actual de esta Política Europea desde una perspectiva funcional, no material. La razón es que, aunque mucho se podría decir sobre el enfoque que adopta presupuestariamente la Política de Cohesión actual, este sería válido únicamente durante el Marco Financiero Plurianual vigente. Por tanto, considero que es más adecuado presentarlo desde una perspectiva conceptual-funcional que permita comprender cómo funciona, cómo se instrumenta, cómo se negocia, independientemente de la estructura financiera vigente. (Que, por otro lado, ya forma parte de otro tema, [\authorfont{Ver Tema 3.B.37}]).

El marco actual tiene sus orígenes en el último periodo presentado, donde se consolida como una Política Integral, multidimensional e imprescindible para alcanzar el resto de objetivos de la Unión.\footnote{%
    Podremos notar que no hemos hecho mención alguna a la última reforma de la norma originaria. La razón es que el Tratado de Lisboa simplemente explicitó la vertiente territorial de la política de cohesión. Es decir, cristalizó una práctica que ya se venía realizando, fruto del periodo previo, y que consiste en buscar un desarrollo equilibrado, prestando una mayor atención a las cuestiones rurales o regiones insulares. Según la Comisión se ha podido constatar que aunque se han alcanzado avances significativos en las regiones menos desarrolladas de la UE en los todos los periodos de programación de la Política de Cohesión, aún persisten importantes diferencias en términos de renta per cápita entre sus regiones.
}


\subsubsection{Dimensión financiera}
Los instrumentos financieros de la Unión Europea se configuran como medidas integradas en el Presupuesto Comunitario con el objetivo de hacer frente a uno o varios objetivos específicos. Tales instrumentos pueden tomar la forma de inversiones, préstamos o garantías, u otros instrumentos de riesgo compartido, y podrán combinarse con subvenciones. La ejecución de la dotación financiera de la Política de Cohesión está en manos de la Comisión. No obstante, alrededor del $70\%$ de los programas se realizan mediante modelos de gestión compartida. Es decir, modelos que implican tanto a la Comisión Europea como a las administraciones nacionales, en lo que respecta a la financiación y la ejecución concreta del programa.

Aunque es difícil definir unos objetivos comunes para esta dotación financiera pues, como veremos, se articula a través de diferentes fondos de naturaleza y objetivos diversos, sí podemos decir que la reducción de las limitaciones estructurales que frenan el progreso de los diversos sectores económicos y alcanzar un mayor grado de cohesión intracomunitaria, exige perseguir, simultaneamentes dos objetivos: 
\begin{la}
    \item \textbf{Inversión en crecimiento y empleo}. Por un lado, el crecimiento económico de las regiones y del empleo en estas. Constituye el objetivo prioritario cuyo objetivo es reducir las disparidades entre las regiones europeas, concentrándose en las regiones y los Estados miembros menos desarrollados, creando condiciones más favorables para el crecimiento y el empleo mediante el aumento de la inversión en capital físico y humano, y la mejora de su calidad, el desarrollo de la innovación y de la sociedad del conocimiento, la lucha contra el desempleo, la protección y mejora del medio ambiente y la eficiencia administrativa; y,
    \item \textbf{Cooperación territorial europea}. Por otro lado, dotar de suficiente apoyo financiero los programas que busquen reforzar la cooperación e integración territorial entre las diferentes regiones europeas. Esto proporciona un marco para el intercambio de experiencias entre agentes locales, regionales y nacionales, intensificando la cooperación transfronteriza a través de iniciativas locales y regionales conjuntas, y fortalece la cooperación transnacional. Así mismo, permite la elaboración de proyectos y programas que, con este fin, puedan ser financiados, total o parcialmente, por la dotación financiera de la que dispone la Comisión, sabiendo que los programas han surgido del consenso y voluntad de los destinatarios; pudiendo aumentar de esta forma las probabilidades del éxito.
\end{la}

Por lo tanto, como vemos la Política de Cohesión es un concepto muy amplio, que engloba multitud de actuaciones diferentes cuyo objeto común es el desarrollo de las regiones y la convergencia real. Esto hace díficil presentar una clasificación uniforme y consensuada sobre los diferentes instrumentos financieros de los que se dispone. Sin embargo, atendiendo a las definiciones jurídicas de los diferentes intrumentos, el conjunto de sus actuaciones y su finalidad, podemos distinguir dos tipos de Fondos principales.\footnote{%
    Esta clasificación deriva de los Reglamentos comunitarios, no de la literatura científica. El Reglamento 1303/2013, por ejemplo, establece expresamente que FEDER y FSE constituyen conjuntamente los \textit{Fondos Estructurales}, mientras que el Fondo de Cohesión tiene un régimen diferenciado. Además, el FEMP (hoy FEMPA) compartía reglas comunes, pero nunca fue considerado un Fondo Estructural.
}

\begin{specialblock}{Marco Financiero Plurianual 2021-2027: Dimensión financiera}     
    En el período 2021-2027, la financiación de la Unión procede de dos fuentes: el marco financiero plurianual (MFP) «clásico», que esboza los límites de los gastos anuales de la Unión, y el plan de recuperación NextGenerationEU, una iniciativa extraordinaria de recuperación diseñada para ayudar a los Estados miembros a recuperarse de la pandemia de COVID-19. La política de cohesión también está financiada por el MFP y, en ciertos casos, por NextGenerationEU. Los recursos destinados al objetivo «Inversión en crecimiento y empleo» de la política de cohesión ascienden a un total de 322 300 millones EUR (a precios de 2018, es decir, en términos de la moneda en 2018) y se asignan como sigue: — 202 300 millones EUR para las regiones menos desarrolladas; — 47 800 millones EUR para las regiones en transición; — 27 200 millones EUR para las regiones más desarrolladas; — 42 600 millones EUR para los Estados miembros que reciben ayudas del Fondo de Cohesión (de los que 10 000 millones EUR se destinarán al Mecanismo «Conectar Europa»); — 1 928 millones EUR en concepto de financiación adicional para las regiones ultraperiféricas, y — 500 millones EUR para inversiones interregionales en innovación. Los recursos procedentes del FEDER correspondientes al objetivo de «cooperación territorial europea» (Interreg) ascienden a un total de 8 050 millones EUR y se asignan como sigue: — 5 812 millones EUR para la cooperación transfronteriza marítima y terrestre; — 1 466 millones EUR para la cooperación transnacional; — 490 millones EUR para la cooperación interregional, y — 281 millones EUR para la cooperación con las regiones ultraperiféricas. El nuevo FTJ, que apoya a los territorios más afectados por la transición hacia la neutralidad climática y tiene por objeto prevenir el aumento de las disparidades regionales, dispone de un presupuesto de 17 500 millones EUR (a precios de 2018, 19 700 millones EUR a precios de 2024). De estos, 7 500 millones EUR proceden del MFP y 10 000 millones EUR adicionales de NextGenerationEU. En diciembre de 2020 se adoptó otro nuevo instrumento, REACT-UE, que sirvió de complemento para los programas de cohesión 2014-2020 y se sumó a las asignaciones de cohesión para el período 2021-2027. REACT-UE prestó apoyo a los sectores más importantes para lograr una recuperación sólida tras la crisis de la COVID-19. Su asignación (hasta 2023) ascendió a 47 500 millones EUR.

    \textbf{Modificaciones de la Política de Cohesión para el periodo: 2021-2027}. El Reglamento sobre disposiciones comunes fue modificado en octubre de 2022 para facilitar y flexibilizar la utilización por parte de los Estados miembros y las regiones de los fondos de la política de cohesión establecidos en el marco de los períodos de programación 2014-2020 y 2021-2027 con el fin de apoyar las medidas destinadas a hacer frente a los retos migratorios resultantes de la agresión militar por parte de la Federación de Rusia contra Ucrania. Para reforzar la competitividad industrial de la Unión mediante el avance en tecnologías críticas, en 2024 se aprobó el Reglamento sobre la Plataforma de Tecnologías Estratégicas para Europa (STEP). STEP apoya las inversiones en proyectos en las fases de desarrollo y fabricación en tres sectores clave para las transiciones ecológica y digital: tecnologías digitales e innovación de tecnología profunda; tecnologías limpias y eficientes en el uso de los recursos, y biotecnologías. A este respecto, las autoridades nacionales pueden reorientar parte de sus fondos de la política de cohesión, como el Fondo de Cohesión, el FEDER, el FSE + y el FTJ, a actividades que apoyen los objetivos de STEP. En diciembre de 2024, se aprobó el Reglamento sobre el Apoyo regional urgente para la reconstrucción (Restore) para proporcionar ayuda adicional a los Estados miembros afectados por catástrofes naturales ocurridas en 2024 y 2025. Este Reglamento permite a los Estados miembros reprogramar fondos del FEDER, el Fondo de Cohesión y el FSE+ para cubrir medidas como la rehabilitación de infraestructuras y equipamientos dañados, el suministro de alimentos y asistencia material básica, y el apoyo sanitario. En septiembre de 2025, tras la revisión intermedia de la política de cohesión, el Parlamento y el Consejo adoptaron el Reglamento (UE) 2025/1914, por el que se modifican los Reglamentos (UE) 2021/1058 (FEDER/Fondo de Cohesión) y (UE) 2021/1056 (FTJ). Las modificaciones permiten a los Estados miembros reasignar los fondos de cohesión de 2025 a las nuevas prioridades estratégicas de la Unión a partir de principios de 2026. Estas prioridades reflejan la evolución del contexto geopolítico y económico y se centran en la competitividad, la defensa y la seguridad, la vivienda asequible, la resiliencia hídrica, la transición energética y el apoyo a las regiones fronterizas orientales. Más concretamente, las modificaciones adoptadas incluyen: — nuevos objetivos del FEDER y del Fondo de Cohesión para las capacidades industriales de defensa y las infraestructuras de doble uso para la movilidad militar y la preparación civil; — apoyo del FEDER a grandes empresas en regiones menos desarrolladas y en transición que contribuyan a la descarbonización o a proyectos importantes de interés común europeo (PIICE) (es decir, proyectos que apoyen, por ejemplo, las infraestructuras de hidrógeno); — aumento de la financiación para viviendas asequibles, con una mayor cofinanciación (hasta el 100 %) y una prefinanciación única del 20 %; — mayor inversión en resiliencia hídrica, infraestructura digital y prevención de sequías; — apoyo a los interconectores energéticos y a las infraestructuras de almacenamiento y recarga; — condiciones preferentes para las regiones fronterizas orientales, entre ellas una prefinanciación adicional de hasta el 9,5 % y una mayor cofinanciación, y — flexibilidad para transferir financiación del FEDER y del Fondo de Cohesión a la Iniciativa Urbana Europea, el Instrumento de Inversiones Interregionales en Innovación o los compartimentos de InvestEU. El Reglamento entró en vigor el 19 de septiembre de 2025.

    Asimismo, el Pacto Verde Europeo pasa a formar parte esencial de la estrategia de recuperación de la UE. Los fondos de cohesión con este objetivo de recuperación avanzan sobre lo acordado en la propuesta de 2018 en:
    - Hincapié en la mejora de la competitividad económica: investigación e innovación, transfonnación digital, Pacto Verde Europeo y refuerzo del pilar de derechos sociales.
    - Flexibilización en la transferencia de recursos de los fondos europeos a los Estados miembros.
    Trabajar en una política de cohesión que sea capaz de responder a situaciones tan excepcionales como la producida por el COVID I 9.
    - Ampliación de plazos para que los Estados miembros completen proyectos que deberían de haber finalizado en el periodo 2014-2020. 
    Reforzar los sistemas de salud y mejorar el potencial de los recursos económicos producidos por la cultura y el turismo. Apoyo al empleo, al empleo joven y a la lucha contra la pobreza infantil.
    
    Durante el período de programación 2021-2027, la Unión asignó más de $392.000\text{M}€$ millones EUR a la política de cohesión. Unos $226.000\text{M}€$ fueron destinados al FEDER.
\end{specialblock}


\paragraph{Fondos Horizontales o Generalista: FSE+ (1958) y FEDR (1975)}
En primer lugar, y con una dotación sustancialmente más importante, tendríamos los Fondos de carácter horizontal (o generalista). Estos Fondos, también conocidos como \textbf{Fondos Estucturales}, financian una enorme variedad de políticas, poseen múltiples objetivos específicos y están presentes en prácticamente todas las regiones; por ello su carácter horizontal. Además, a fin de garantizar un uso eficiente de sus recursos, comparten y deben respetar una serie de principios comunes: (i) \textbf{concentración}, las intervenciones deben estar principalmente destinadas a financiar proyectos en las regiones menos avanzadas; (ii) \textbf{cooperación}, entre la Comisión y las autoridades de los Estados miembros (nacionales regionales y locales); (iii) \textbf{programación}, programación plurianual de las intervenciones a través de la elaboración de programas que se ejecutarán durante varios años; y, (iv) \textbf{adicionalidad}, las aportaciones comunitarias no deben sustituir a los gastos nacionales, deben ser complementarios (cofinanciación).  Asimismo, y como muestra de su carácter generalista, para la gestión de los Fondos \textbf{se adopta entonces la nomenclatura común de las Unidades Territoriales Estadísticas (NUTS)}, incorporadas por Eurostat con el fin de ofrecer una división uniforme de unidades territoriales para la elaboración de las estadísticas regionales de la Unión Europea. Los Fondos Estructurales son dos: el Fondo Social Europeo y el Fondo Europeo de Desarrollo Regional.

\textbf{Fondo Social Europeo}. El Fondo Social Europeo se crea en 1958 con el fin de permitir la creación de empleo, mejorar la formación y cualificación de la mano de obra, facilitar la integración de los trabajadores del sur de Europa y su desplazamiento. Este objetivo respondía a las necesidades del mercado de trabajo de la época, caracterizado por un exceso de demanda en las zonas con más dinamismo de importantes flujos migratorios debido a la expansión económica en la que estaba inmersa la economía mundial durante estos años. Llamado FSE+ desde 2021, es el principal instrumento de la Unión por lo que respecta al apoyo de medidas encaminadas a prevenir y combatir el desempleo, desarrollar los recursos humanos y fomentar la integración social en el mercado de trabajo. Financia iniciativas que promueven un nivel elevado de empleo, la igualdad de oportunidades para mujeres y hombres, el desarrollo sostenible y la cohesión económica y social.

\textbf{Fondo Europeo de Desarrollo Regional}. Por otro lado, estaría el Fondo Europeo de Desarrollo Regional, creado en 1975, cuyo fundamento jurídico se encuentra en los arts. 174 a 178 del TFUE. El FEDER fue y es un Fondo destinado a corregir los principales desequilibrios regionales en la CEE de entonces y en la actual UE. Es uno de los principales instrumentos financieros de la política de cohesión europea (económica, social y territorial). Su objetivo es contribuir a reducir las diferencias entre los niveles de desarrollo de las regiones europeas y el retraso de las regiones menos favorecidas, fortaleciendo la cohesión económica, social y territorial en la Unión Europea corrigiendo los desequilibrios entre sus regiones, y apoyando la inversión en crecimiento y empleo. De hecho, el artículo 176 del TFUE dispone que el FEDER está destinado a contribuir a la corrección de los principales desequilibrios regionales dentro de la Unión, y para lograrlo conceder apoyo al desarrollo y al ajuste estructural de las regiones menos desarrolladas y a la reconversión de las regiones industriales en declive. Como veremos, los recursos del FEDR se ejecutan mediante programas de gestión compartida y, en consecuencia, juegan un papel fundamental los Acuerdos de Asociación que luego veremos, donde se fijan los pormenores de la asignación y la utilización de los fondos del FEDER. El marco general para la ejecución y puesta en práctica del FEDER se establece en el Reglamento sobre Disposiciones Comunes, que abarca otros fondos de la Unión y determina objetivos específicos y el ámbito de aplicación de las medidas de apoyo posibles.






















La Comisión mide la riqueza a través del Producto Interior Bruto (PIB) y la Renta Nacional Bruta (RNB) para reflejar el poder adquisitivo. Esta categorización es uno de los aspectos que se ha visto más afectado por el período de programación del MFP 202 1-2027 ya que se modula este criterio de reparto con nuevos criterios (desempleo juvenil, bajo nivel educativo, cambio climático y recepción e integración de inmigrantes) para reflejar mejor la realidad de cada una de las zonas, de tal manera que alguna región que en el período anterior 20 14-2020 perteneciera a la categoría de regiones más desarrolladas, puede pasar a la categoría de regiones en transición en el siguiente período. El efecto de pasar de un estatus de «más desarrollada» a «en transición» es que la región en cuestión tiende a recibir más financiación del ámbito de cohesión.




\begin{specialblock}{Marco Financiero Plurianual 2021-2027: Fondo Europeo de Desarrollo Regional}
    Las normas para el FEDER en el período 2021-2027 se establecen en: el \href{https://www.boe.es/buscar/doc.php?id=DOUE-L-2021-80891}{Reglamento relativo al FEDER y al Fondo de Cohesión} y el \href{https://www.boe.es/buscar/doc.php?id=DOUE-L-2021-80892}{Reglamento sobre disposiciones específicas para el objetivo de cooperación territorial europea (Interreg)}. En este MFP se han establecido cinco objetivos que deben guiar las acciones de los Fondos Estructurales (también el FC, que veremos después):\footnote{%
        El FEDER también respalda el desarrollo urbano sostenible. En el período 2021-2027, al menos el $8\%$ de los recursos del FEDER (a nivel nacional) se destinan al desarrollo urbano sostenible y a la creación de la Iniciativa Urbana Europea, que permite a las áreas urbanas experimentar con soluciones innovadoras para hacer frente a los retos de las zonas urbanas. La política de cohesión cuenta con una lista de actividades que no reciben fondos del FEDER, como el desmantelamiento o la construcción de centrales nucleares, infraestructuras aeroportuarias (excepto en las regiones ultraperiféricas), algunas operaciones de gestión de residuos (como los vertederos) y el apoyo a la industria del tabaco.
    }
    \begin{la}
        \item \textbf{Objetivo I: Europa Inteligente}. Una Europa más competitiva e inteligente, promoviendo una transformación económica innovadora e inteligente y una conectividad regional a las tecnologías de la información y de las comunicaciones, mediante la innovación, la digitalización, la transformación económica y el apoyo a las pequeñas y medianas empresas;
        \item \textbf{Objetivo II: Europea más verde}. Una Europa más verde, baja en carbono, en transición hacia una economía climáticamente neutra y resiliente, promoviendo una transición energética limpia y justa, la inversión verde y azul, la economía circular, la mitigación y adaptación al cambio climático, la prevención y gestión de riesgos y la movilidad urbana sostenible, que aplique el Acuerdo de París e invierta en transición energética, energías renovables y la lucha contra el cambio climático;
        \item \textbf{Objetivo III: Europea conectada}. Una Europa más conectada, mejorando la movilidad, con un transporte estratégico y redes digitales;
        \item \textbf{Objetivo IV: Europa social}. Una Europa más social e intuitiva, por medio de la aplicación del pilar europeo de derechos sociales, que haga realidad el pilar europeo de derechos sociales y que apoye el empleo de calidad, la educación, las capacidades educativas y profesionales, la inclusión social y la igualdad de acceso a la asistencia sanitaria; y,
        \item \textbf{Objetivo V: Europa y sus ciudadanos}. Una Europa más cercana a los ciudadanos, que respalde estrategias de crecimiento de gestión local y que contribuya a un desarrollo urbano sostenible en toda la UE, fomentando el desarrollo integrado y sostenible de todo tipo de territorios e iniciativas sociales.
    \end{la}
    
    Las estipulaciones respecto a las asignaciones que debe tener cada programa no son muy abundantes. Sin embargo, sí resulta obligatorio destinar al menos el $25\%$ (del FEDER) al Objetivo 2 y concentrar el gasto en el Objetivo 1, aunque, claro está, atendiendo a los diferentes niveles de properidad y necesidades.\footnote{%
        Más aún, el Reglamento establece que:
    \begin{la}
        \item Las regiones o los Estados miembros menos desarrollados tienen que dedicar al menos el $25\%$ al Objetivo 1;
        \item Las regiones o Estados miembros en transición tienen que dedicar al menos el $40\%$ al Objetivo 1;
        \item Las regiones o los Estados miembros más desarrollados tienen que dedicar al menos el $85\%$ de su asignación al Objetivo 1 y al Objetivo 2.
    \end{la}
    } Por otro lado, en lo relativo al régimen de cofinanciación, el Reglamento estipulata que los proyectos financiados mediante la asignación de recursos provenientes de Fondos Estructurales atenderá lo siguiente:
    \begin{la}
        \item \textbf{Regiones menos desarrolladas: $\leq 85\%$}. Los programas realizados en las regiones menos desarrolladas, aquellas cuyo PIBc sea $\leq75\%$ de la media de la UE, podrán cofinanciarse hasta el $85\%$. En España, incluiría las Comunidades Autónomas de Andalucía, Castilla La Mancha, Ceuta, Extremadura y Melilla;
        \item \textbf{Regiones en transición: $\leq 60\%$}. Los programas realizados en las regiones en transición, aquellas cuyo PIBc sea $100>x\geq75\%$ de la media de la UE, podrán cofinanciarse hasta el $60\%$. En España, incluiría las Comunidades Autónomas de Asturias, Baleares, Canarias, Cantabria, Castilla León, Galicia, La Rioja, Murcia, Valencia; y,
        \item \textbf{Regiones avanzadas: $\leq 40\%$}. Los programas realizados en las regiones más avanzadas, aquellas cuyo PIBc sea $\geq100\%$ de la media de la UE, podrán cofinanciarse hasta el $50\%$. En España, incluiría las Comunidades Autónomas de Aragón, Cataluña, Navarra, Madrid, País Vasco.
    \end{la}

    Durante el período de programación 2021-2027, la Unión asignó unos $226.000\text{M}€$ al FEDER.
\end{specialblock}

\paragraph{Fondos de Inversión: FC (1993), FEMP (2014) y FTJ (2019)}
Por otro lado estarían los Fondos Finalistas o Especializados, aunuque también se conocen en el ámbito comunitario como los Fondos de Inversión. Estos Fondos se caracterizan por tener un ámbito material claramente delimitado,  sólo financian determinadas actuaciones y estar dirigidos a determinadas regiones/sectores, que son los únicos potenciales beneficiarios.

Esta categoría de Fondos abarca un conjunto de instrumentos amplío y extremedamente heterogeneo. Sin ánimo de ser exhaustivo, cabría destacar el Fondo de Cohesión, el Fondo Europeo Martítimo y de Pesca y el Fondo de Transición Justa. 

\textbf{Fondo de Cohesión}. El Fondo de Cohesión inicia su andadura en 1993, mediante el Tratado de Maastricht, con la misión de incrementar la dotación financiera de los países miembros con un PIB inferior al $90\%$ de la media (CEE$=12$). Inicialmente, sus ayudas estaban dirigidas a financiar proyectos para mejorar el medio ambiente y financiar las infraestructuras de transporte de acuerdo con lo estipulado en la regulación de las redes transeuropeas de transporte (especialmente la alta velocidad europea, RTE-T). Hoy en día, su fundamento jurídico se encuentra en el art. 177, p.2º, del TFUE, que establece el Fondo de Cohesión en consonancia con los objetivos de cohesión económica, social y territorial de la Unión Europea. A lo largo de los años, ha servido para impulsar la política de cohesión de la Unión Europea, tanto a través de la financiación de redes de transporte como tambiém, desde 2007, del apoyo a sectores vinculados al desarrollo sostenible. La asignación de sus recursos mantiene criterios de renta, por lo que su actuación no está presente en todas las regiones. De hecho, España no ha sido elegible desde 2014.\footnote{%
    No obstante ha seguido recibiendo, en los primeros años de dicho periodo, ingresos procedentes de la ayuda transitoria del régimen de \textit{phasing out} del Fondo de Cohesión en el Objetivo de Convergencia al que tuvo acceso en el periodo 2007-2013, junto con pequeños saldos correspondientes al cierre de proyectos pertenecientes al periodo 2000-2006. En el ejercicio 2018 se recibió en España el último importe correspondiente al Fondo de Cohesión.
} 

Su ámbito de actuación es amplío, apoyando inversiones en medio ambiente, desarrollo sostenible, energía (en particular las energías renovables), transporte, asistencia técnica, información, comunicación y estudios. La manera de actuar suele ser compartida, donde el Fondo de Cohesión financia programas en los que la responsabilidad recae tanto en la Comisión Europea como las autoridades nacionales y regionales. Lo más habitual, como se ha dicho, es que las autoridades locales elaboren sus propuestas, soliciten financiación y, en consecuencia, asuman responsabilidad en su gestión cotidiana, mientras que la Comisión contribuye a garantizar que los proyectos concluyen con éxito y comprueba que los fondos se utilizan correctamente.

\textbf{Fondo de Transición Justa}. Por otro lado tendríamos el Fondo de Transición Justa. Este Fond es fruto del conocico \textbf{Pacto Verde Europeo}, adoptado en 2019 por la Comisión, que establece una hoja de ruta para una nueva política de crecimiento en la Unión, más sostenible y con el objetivo de alcanzar la neutralidad climática. En este sentido, el FTJ se constituye como uno de los tres pilares principales del pacto (junto al marco InvestEU y una nueva estrategia para el BEI), cuyo objetivo principal es lograr la \textbf{diversificación económica de los territorios más afectados por la transición climática}, así como en el reciclaje y la inclusión activa de sus trabajadores y demandantes de empleo. Para este fin concede principalmente subvenciones cuyo acceso requiere cumplir unos criterios de admisibilidad. No obstante, estos son más amplios a fin de apoyar también actividades relacionadas con la transición energética. De hecho, normalmente se distancia de los criterios de renta y parte de criterios basados en las emisiones industriales, el empleo en la industria y en la minería de combustibles fósiles, y el nivel de desarrollo económico. 

Se ejecuta mediante reglas de gestión compartida, lo cual implica una estrecha cooperación con las autoridades nacionales, regionales y locales. A fin de acceder al apoyo, los Estados miembros tienen que presentar planes territoriales de transición justa.\footnote{%
    De hecho, aquellos Estados miembros no comprometidos a alcanzar la neutralidad climática en 2050, recibirán exclusivamente el $50\%$ de su asignación prevista.
} Estos planes indicarán las áreas concretas de intervención, basándose en las repercusiones económicas y sociales derivadas de la transición. En particular, estos planes deben tener en cuenta las pérdidas de puestos de trabajo previstas y la transformación de los procesos de producción de las instalaciones industriales con la mayor intensidad de gases de efecto invernadero. \textbf{Proporciona apoyo a todos los Estados miembros}, pero el nivel de cofinanciación se fija en función de la categoría de la región en la que se hallen estos proyectos.\footnote{%
    Para las regiones menos desarrolladas, se fija en un máximo del $85\%$, para las regiones en transición, en un $70\%$ y para las regiones más desarrolladas, en un $50\%$.
} 

En definitiva, es una herramienta clave para apoyar a los territorios más afectados por la transición a la neutralidad climática y evitar el aumento de las disparidades regionales. 

\begin{specialblock}{Marco Financiero Plurianual 2021-2027: Fondo Especializados}
    \textbf{Fondo de Cohesión}. En el período comprendido entre 2021 y 2027, el Fondo de Cohesión financia: inversiones en medio ambiente, incluidos los ámbitos relacionados con el desarrollo sostenible y la energía que presenten beneficios para el medio ambiente, redes transeuropeas en materia de infraestructuras de transporte (por ejemplo, la RTE-T) y asistencia técnica.\footnote{%
        La política de cohesión establece una lista de actividades que no pueden recibir apoyo del Fondo de Cohesión en el período 2021-2027, como el desmantelamiento o la construcción de centrales nucleares, las infraestructuras aeroportuarias (excepto en las regiones ultraperiféricas) y algunas instalaciones de gestión de residuos (como los vertederos).
    } Sus normas se han fijado en el Reglamento relativo al Fondo Europeo de Desarrollo Regional y al Fondo de Cohesión. De entre los objetivos establecidos para los Fondos Estructurales, el Fondo de Cohesión persigue únicamente dos objetivos: una economía verde (Objetivo II) y una Europa más conectada (Objetivo III).

    El presupuesto de la Unión para el Fondo de Cohesión en el período 2021-2027 está fijado en $42.600\text{M}€$, incluidos $10.000\text{M}€$ millones EUR para el Mecanismo \textbf{Conectar Europa}, un programa de financiación de la Unión cuyo objetivo es impulsar el desarrollo de una infraestructura transeuropea en ámbitos como el transporte, la energía y los servicios digitales. El porcentaje máximo de cofinanciación puede llegar al $85\%$ del valor de los proyectos. Se espera que el $37\%$ del total de las asignaciones financieras del Fondo de Cohesión contribuya a los objetivos climáticos de la Unión.

    En el período, ha proporcionado apoyo a quince países de la Unión: Bulgaria, Chequia, Chipre, Croacia, Eslovaquia, Eslovenia, Estonia, Grecia, Hungría, Letonia, Lituania, Malta, Polonia, Portugal y Rumanía.

    \textbf{Fondo de Transición Justa}. El Fondo de Transición Justa cuenta con un presupuesto global de $17.500\text{M}€$ para el período 2021-2027. Los Estados miembros pueden complementar su asignación del Fondo de Transición Justa con los recursos que perciban en virtud del Fondo Europeo de Desarrollo Regional y el Fondo Social Europeo Plus.
\end{specialblock}

\begin{specialblock}{Otros Intrumentos Financieros de carácter finalista}
    Existen otros instrumentos financieros de carácter finalista que no es necesario tratar. En particular destacan:
    \begin{la}
        \item \textbf{Reserva de Adaptación al Brexit}. Esta Reserva va dirigida a proporcionar apoyo a los Estados miembros, las regiones y los sectores más afectados por las consecuencias adversas de la retirada del Reino Unido de la Unión; y en particular contribuirá a las medidas específicas establecidas por los Estados miembros para ayudar a las empresas y los sectores económicos, los trabajadores, las regiones y las comunidades locales que sufren el impacto del final del período de transición. En concreto. la contribución financiera de la Reserva solo apoyará el gasto público directamente relacionado con las medidas adoptadas específicamente por los Estados miembros para contribuir a los objetivos establecidos en el correspondiente \href{https://eur-lex.europa.eu/eli/reg/2021/1755/oj/spa}{Reglamento}. En cuanto a la asignación, la totalidad asciende a 5.000 millones de euros (a precios de 2018), que se asignarán provisionalmente por adelantado a los Estados miembros; y,
        \item \textbf{Acción de Cohesión para los Refugiados en Europa (CARE)}. Creado por el \href{https://www.boe.es/buscar/doc.php?id=DOUE-L-2022-80558}{Reglamento 2022/562}, permite a los Estados miembros y las regiones prestar asistencia urgente a las personas que huyen de la invasión de Ucrania por parte de Rusia. CARE introduce la flexibilidad necesaria en las normas de la política de cohesión 2014-2020 para permitir una rápida reasignación de la financiación disponible para esa asistencia urgente. Además, la dotación para 2022 de los fondos de Ayuda a la Recuperación para la Cohesión y los Territorios de Europa (REACT-UE), de 10 000 millones de euros, también puede utilizarse para hacer frente a estas nuevas necesidades, en el marco del objetivo general de la recuperación posterior a la pandemia. CARE introduce cuatro camoios principales en las normas de la política de cohesión para maximizar la rapidez y facilidad con la que los Estados miembros pueden ayudar a las personas que huyen de Ucrania, manteniendo al mismo tiempo las ayudas para la recuperación de las regiones.
    \end{la}
\end{specialblock}

\subsubsection{Dimensión política}
Al margen de esta actividad financiera, la Política de Cohesión se articula también a través de una dimensión política. 

Esta dimensión política, muy relevante, permite identificar, planificar y coordinar las actividades de la Unión con las actividades de las autoritaridades nacionales, bien sean de carácter estatal o de carácter regional/local. Se articula, fundamentalmente, a través de dos instrumentos/formas principales: los Acuerdos de Asociación, de carácter jurídico y estatal, y el Comité Europeo de las Regiones, foro formal de carácter regional.

\paragraph{Planificación estratégica: Acuerdo de Asociación}
El Acuerdo de Asociación (AA) es un documento de carácter estratégico, elaborado por cada uno de los Estados miembros que recoge las prioridades nacionales de inversión y el planteamiento básico de las actuaciones que se llevaran a cabo en materia de cohesión territorial con los recursos de los que se dispone a través de los diferentes Fondos (Estructurales y de Inversión). 

Estos Acuerdos son inicialmente una propuesta elaborada, normalmente, con la participación de socios regionales (entre otros, Gobiernos regionales y locales) y socios sociales (por ejemplo, organizaciones de la sociedad civil y sindicatos). En el documento se elabora la estrategia de desarrollo del estado, que identifica sus necesidades y prioridades de inversión, los objetivos y los resultados esperados de la financiación europea. Además, establece los procedimientos de gestión y control de la financiación europea, así como los mecanismos de seguimiento y evaluación de los resultados. El objetivo es garantizar que los fondos se utilicen de manera eficaz y eficiente para lograr los objetivos de la estrategia de desarrollo y cumplir con las normas y reglamentos de la UE. Una vez elaborados son enviados a la Comisión y si ésta considera que las actividades, programas y objetivos se alinean con las directrices fijadas para la asignación de estos Fondos a nivel comunitario, la propuesta deviene un acuerdo entre la Comisión y el Estado miembro sobre cómo utilizar los recursos. 

Dentro de estas propuestas/acuerdos es abitual encontrar una estructuración por \textbf{Programas Operativos} (PO). Los programas operativos son planes detallados en los que se establece cómo se gastará el dinero procedente de los Fondos durante el período de programación, en linea con la Nueva Teoría de la Gestión Público basada en el presupuesto orientado a resultados y la planificación estratégica plurianual.  Los programas operativos pueden elaborarse: (i) para una región específica o para un objetivo temático de alcance nacional (por ejemplo, medio ambiente); o, (ii) alrededor de varios Programas Operativos (PO) plurirregionales (nivel nacional, como el POPE).\footnote{%
    \textbf{NOTA.-} Es interesante destacar que, en caso de Programas Operativos plurriregionales es posible que se observe la participación de regiones de diferentes Estados Miemrbos, cumpliendo el objetivo de Cooperación Territorial Europea. Estos programas operativos interregionales o transfronterizos que afectan a diferentes Estados Miembros. Además, cabe destacar que, en el caso de España, aunque la AGE es el receptor de los Fondos, existe un reparto interno a $50\%$ entre AGE y Administraciones Autonómicas.
} Además, estos Programas Operativos especificarán quien actúa como:
\begin{la}
    \item \textbf{Autoridad de Gestión (AG)}. Es una autoridad pública, nacional, regional o local o un organismo público o privado designado por el estado para gestionar y ejecutar un programa operativo. Realiza funciones relativas a la gestión del programa operativo, la selección de las operaciones o la gestión y control financieros del programa operativo. Es la responsable de la gestión y aplicación efectiva de un programa operativo;
    \item \textbf{Autoridad de Auditoría (AA)}. Es una autoridad u organismo público, nacional, regional o local, funcionalmente independiente de las autoridades de gestión y de certificación, designada por el estado y responsable de verificar el funcionamiento efectivo del sistema de gestión y control. Entre las tareas de la autoridad de auditoría se incluye la comprobación del funcionamiento eficiente de los sistemas de control y gestión (auditorías de los sistemas). También es responsable de llevar a cabo los controles destinados específicamente al gasto declarado (auditorías operativas). En España, las funciones de la autoridad de auditoría las asume la Intervención General del Estado (IGE);
    \item \textbf{Autoridad de Certificación (AC)}. Es una autoridad u organismo público, nacional, regional o local designado por el estado a fin de garantizar la precisión de las declaraciones de gastos y las solicitudes de pago antes de su envío a la Comisión Europea. Por tanto, es la responsable de comprobar, por un lado, que la declaración de gastos es exacta, que se ha realizado aplicando sistemas de contabilidad fiables y que se basa en justificantes verificables y, por otro, que el gasto declarado se atiene al derecho aplicable en la materia y se ha realizado en relación con las operaciones seleccionadas para financiación, de confcmnidad con los criterios aplicables al programa y en cumplimiento del derecho aplicable; 
    \item \textbf{Organismos Intermedios}. Se definen como todo organismo público o privado que, bien actúa bajo la responsabilidad de una autoridad de gestión o de certificación para la realización de determinadas tareas propias de esas autoridades, bien desempeña competencias en nombre de tal autoridad en relación con los beneficiarios que ejecuten las operaciones en las que se concreta la ejecución de un programa operativo determinado; y,
    \item \textbf{Organismo al que la Comisión Europea debe hacer los pagos}. En España es el Ministerio de Hacienda y Función Pública quien recibe los Fondos y los reparte en función de los distintos Programas Operativos a través de la Secretaría General de Fondos Europeos.
\end{la}

Asimismo, cabe destacar que la Comisión puede solicitar a un Estado miembro modificar su Acuerdo de Asociación y los programas considerados relevantes para contribuir a la aplicación de las Country Specific Recommendations (CSR), pudiendo suspender los pagos al Estado miembro si éste no adopta medidas eficaces.\footnote{%
    \textbf{Country Specific Recommendations (CSR)}. La posibilidad que acabamos de comentar tiene su razón de ser en garantizar la coherencia de las Políticas de la Unión con las observaciones y resultados con el Semestre Europeo. El Semestre Europeo es un foro que no analizaremos en este Tema. Brevemente, es un ciclo de coordinación de las Políticas Económicas (presupuestarias, sociales y de empleo) dentro de la UE y forma parte del marco de gobernanza económica de la Unión Europea. Además, proporciona un marco de coordinación y supervisión de los esfuerzos de los Estados miembros por cumplir los principios y derechos establecidos en el \textit{pilar europeo de derechos sociales}. Para lo que nos ocupa, lo único que debemos saber es que, el Semestre Europeo sirve de marco para orientar las reformas económicas y sociales que se consideran necesarias en los Estados miembros, así como para hacer un seguimiento de las mismas. De esta forma, al preparar los Acuerdos de Asociación, los Estados miembros tienen que buscar la congruencia en el uso de los Fondos y los Programas Nacionales de Reformao y Country Specific Recommendatios que puedan haber surgido del Foro, en caso de haberlas.
} Además, también podrá suspender los pagos si el Consejo concluye que no está tomando medidas eficaces para corregir un \textit{desequilibrio excesivo} (déficit).

\begin{specialblock}{Marco Financiero Plurianual 2021-2027: Acuerdo de Asociación de España}
    El Acuerdo de Asociación 202 1 -2027 se ha realizado bajo el principio de asociación y gobernanza en varios niveles y supone una inversión total de 59. 722 millones de euros. La suma de los fondos incluidos en el Acuerdo supone un incremento del 15\% respecto a los obtenidos en el periodo2014-2020.

    \textcolor{red}{Nota de Prensa Consejo de Ministros (11 de octubre 2022): Resumen Fondos 2021-2027}

    Esta inversión se organiza en torno a cinco grandes áreas políticas, en las que cobran especial relevancia el apoyo a la investigación, la digitalización, el apoyo a las pymes, la eficiencia energética, la transición verde y la inversión de carácter social.

    Sobre la base de las líneas estratégicas definidas en el Acuerdo, se prevé próximamente la aprobación de los Programas financiados con estos fondos. con la concreción de las tipologías de actuación financiables, de acuerdo con la siguiente estructura: FEDER, contará con una dotación total de 23.397 millones de euros distribuidos en 19 programas regionales y un programa plurirregional. FSE, contará con 11.296 millones de euros, repartidos en 19 programas regionales y 4 programas plurirregionales: Programa Plurirregional de Empleo, Formación Profesional y Educación de España, Programa Plurirregional de empleo Juvenil; Programa Plurirregional de Inclusión y Lucha contra la pobreza y la pobreza infantil, y el Programa Plurirregional de Privación Material. Por su parte, el FEMPA, con 1.120 millones, y el FTJ, con 869 millones, contarán con un solo programa plurirregional cada uno de ellos. 
    
    Los datos se resumen en el siguiente gráfico:\footnote{%
        Asignación financiera por temática y por fondo: Brochure Acuerdo de Asociación 2021-2027
    }
    
    Todos los puntos del Acuerdo de Asociación desembocan en una serie de programas de inversión específicos, que canalizan la financiación a las diferentes regiones y proyectos en los ámbitos de actuación correspondientes. Se denominan Programas Operativos. En cada programa operativo el Estado determina cuál de los objetivos temáticos (en los que se concentra la política de cohesión para el período de programación en marcha) se abordarán con los fondos disponibles para dichos programas. La UE distribuye los fondos a los Estados miembros y regiones en función de lo acordado en cada programa operativo.

    En este sentido, se destaca que España ha sido un receptor histórico de fondos europeos. España fue por primera vez un contribuyente neto entre 20 16 y 2020. Sin embargo, desde 2021. vuelve a ser receptor neto, condición que se ha mantenido en 2022. Respecto al presupuesto del Marco Financiero Plurianual 2021-2027. se calcula que España recibirá durante este tiempo 79.000 millones de euros que incluye 21.000 millones en FEDER y 10.000 millones en FSE+. No obstante, desde el MFP 20 1 3-2020, España ya no recibe Fondos de Cohesión.
\end{specialblock}


\paragraph{Comité de las regiones}
El segundo componente de esta dimensión política sería el Comité Europeo de las Regiones.\footnote{%
    Compuesto por 329 miembros que representan a las entidades regionales y locales de los 27 Estados miembros de la Unión Europea. El CDR emite dictámenes en los casos de consulta obligatoria contemplados en los Tratados, en los casos de consulta facultativa y por propia iniciativa si lo considera oportuno. Sus miembros no están vinculados por ningún mandato imperativo. Ejercen sus funciones con plena independencia, en interés general de la Unión. Su fundamento jurídico se encuentra en el art. 13, apartado 4, del Tratado de la Unión Europea (TUE); artículos 300 y 305 a 307 del Tratado de Funcionamiento de la Unión Europea (TFUE), y varias decisiones del Consejo por las que se nombran miembros y suplentes del Comité Europeo de las Regiones, de conformidad con las propuestas de los Estados miembros, para un mandato de cinco años.
}

Fue creado en 1994 a raíz de la entrada en vigor del Tratado de Maastricht, el Comité Europeo de las Regiones (CDR) es un órgano consultivo de representación de las entidades regionales y locales de la Unión. Por decirlo sencillamente: \textbf{es el portavoz de los intereses de estas entidades ante el Parlamento Europeo, el Consejo y la Comisión, a los que remite dictámenes}.

Sus miembros pueden ser, por ejemplo, altos cargos de entidades regionales, alcaldes o representantes electos o no electos de regiones y ciudades de los veintisiete Estados miembros. Según su propia declaración de misión, el CDR es una asamblea política compuesta de miembros electos, regionales o locales, al servicio de la causa de la integración europea, y garantiza la representación institucional del conjunto de los territorios, regiones, ciudades y municipios de la Unión. Por tanto, busca implicar a las entidades regionales y locales en el proceso europeo de toma de decisiones, facilitando de esta forma la participación de los ciudadanos.

Se trata de un órgano político que aúna y empodera a los representantes locales y regionales elegidos en Europa, con 329 miembros titulares y 329 suplentes, que se escogen entre un millón de políticos regionales y locales. En su conjunto, representan a las $300$ regiones de la Unión (NUTS), a $100000$ entes locales y a cerca de $450$ millones de ciudadanos.

Según sus tres prioridades políticas para 2025-2030, adoptadas el 15 de mayo de 2025, todas las decisiones tomadas a escala de la Unión para abordar las grandes transformaciones sociales a las que se enfrentan actualmente pueblos, ciudades y regiones deben adoptarse lo más cerca posible a los ciudadanos en consonancia con el principio de subsidiariedad:
\begin{la}
    \item \textbf{Prioridad I: Cohesión}. Mejorar la competitividad mediante políticas de base local y ecosistemas de innovación locales; elaborar el presupuesto de la Unión para el período posterior a 2027 en coordinación con los entes locales y regionales; garantizar unas transiciones justas para gestionar las repercusiones socioeconómicas de la triple transición (ecológica, digital y demográfica); y abordar la grave crisis de la vivienda a fin de garantizar una vivienda sostenible y asequible para todos;
    \item \textbf{Prioridad II: Resiliencia}. Adquirir la capacidad de dar respuesta a los desastres naturales y antropogénicos, las amenazas para la seguridad y la inestabilidad económica y política; reforzar la cooperación territorial europea; luchar contra el cambio climático mediante la mitigación y la adaptación; garantizar la resiliencia hídrica; lograr una energía limpia, segura y asequible; establecer una Unión Europea de la Salud; apoyar una política agrícola común sólida y mejorar la preparación y las capacidades de defensa de Europa, también mediante la Alianza para la Reconstrucción de Ucrania; y,
    \item \textbf{Prioridad III: Proximidad}. Acercar Europa a la ciudadanía reforzando la capacidad de los Gobiernos locales; abordar el cambio demográfico mediante la política de cohesión y soluciones de base local; salvaguardar y reforzar la democracia local y la participación de la ciudadanía; incorporar la igualdad de género en todas las políticas de la Unión y garantizar la representación equitativa en la vida política; y aplicar la Estrategia de la UE para la Juventud con vistas a la mejora de las oportunidades de aprendizaje y empleo de los jóvenes.
\end{la}

De entre sus atribuciones más importantes, destaca por encima de cualquier otra su condición de órgano consultivo, de carácter obligatorio para cualquier cuestión relacionado con: (i) educación, formación profesional y juventud (artículo 165 del TFUE); (ii) cultura (artículo 167 del TFUE); (iii) salud pública (artículo 168 del TFUE); (iv) \textbf{redes transeuropeas} de transportes, telecomunicaciones y energía (artículo 172 del TFUE); y, (iv) \textbf{cohesión} económica y social (artículos 175, 177 y 178 del TFUE). Es decir, el Consejo y la Comisión están obligados a solicitar el dictamen del CDR antes de adoptar una decisión relacionada con estas materias, lo que les otorga un papel fundamental en la elaboración, planificación y decisión de la Política de Cohesión.


\subsubsection{Dimensión política}







Por otro lado, la Cooperación Territorial Europea (en adelante, CTE) es el objetivo de la política de cohesión concebido para solucionar aquellos problemas que trascienden las fronteras nacionales y requieren una solución común, y para desarrollar de forma conjunta el potencial de los distintos territorios.

Los programas ejecutados en el marco del objetivo de la CTE, también denominados Interreg, han formado parte de la política de cohesión desde 1990. Las acciones de la CTE reciben apoyo del FEDER. Es el programa insignia para la cooperación transfronteriza a nivel regional y nacional.

Como características reseñables: todas las fronteras terrestres internas y externas de la UE, así como las fronteras marítimas (con condiciones) se enmarcan en la cooperación transfronteriza la Comisión definirá tas zonas correspondientes a la cooperación transnacional y con la opción de que los Estados miembros puedan añadir territorios adyacentes - la cooperación interregional abarcará la totalidad del territorio de la Unión las regiones ultraperiféricas podrán combinar acciones de cooperación transfronteriza y transnacional en un único programa de cooperación. en los programas de cooperación también podrán participar terceros países.

Para el período de programación 202 1 -2027, la CTE para tiene una dotación de aproximadamente 8.000 millones €. Este programa, además de contribuir a los objetivos establecidos para los FEDER, servirá además a dos objetivos específicos: una mejor gobernanza de la cooperación y una Europa más segura y protegida.\footnote{%
    \href{https://www.europarl.europa.eu/factsheets/es/sheet/98/la-cooperacion-territorial-europea}{\textit{La Cooperación territorial europea}. Sede del Parlamento Europeo (Fact Sheet)}
}


\subsection{Covergencia real: ¿Cuáles han sido sus resultados?}
Con cada período presupuestario, la política de cohesión se adapta a los nuevos retos globales y regionales, y continúa siendo un pilar clave para la integración y el desarrollo de la Unión Europea. El octavo informe sobre la cohesión publicado por la Comísión pone de manifiesto que la política de cohesión ha ayudado a reducir las disparidades territoriales y sociales entre las regiones de la U E.\footnote{%
    \textit{Octavo Informe sobre la Cohesión económica, social y territorial. La cohesión en Europa en el horizonte 2050} (Diario Oficial de la Unión Europea). 2022
} Gracias a los fondos de cohesión, se prevé que el PIB per cápita de las regiones menos desarrolladas aumente hasta un 5\% de aquí a 2023. Estas mismas inversiones también contribuyeron a reducir un 3,5\% la brecha entre el PIB per cápita del 10\% de las regiones menos desarrolladas y del 10\% de las regiones más desarrolladas.

Una vez analizada la Política Regional de la Unión Europea, en términos generales, se puede decir que ha habido una convergencia real entre las regiones de la UE en los últimos años. Según los datos de la Comisión Europea, la brecha económica entre las regiones más ricas y las más pobres de la UE se ha reducido en un 34\% desde 2000.\footnote{%
    \href{https://ec.europa.eu/regional_policy/information-sources/publications_es}{PROPUESTA DE RESOLUCIÓN DEL PARLAMENTO EUROPEO sobre la cohesión económica, social y territorial en la UE: octavo infonne sobre la cohesión (2022 '20321 INI ¡)}
} Además, el PIB per cápita de las regiones más pobres ha aumentado a un ritmo más rápido que el de las regiones más ricas.

Sin embargo, esta convergencia no ha sido uniforme en todas las regiones de la UE. Aunque la mayoría de las regiones han mejorado su situación económica, todavía existen diferencias significativas entre las regiones más ricas y las más pobres. Las regiones del sur y del este de Europa, en particular, todavía tienen un nivel de desarrollo económico y social inferior al de las regiones del norte y del oeste de Europa.

En cuanto al papel que ha desempeñado la política regional en la convergencia, se puede decir que ha sido un factor importante. La financiación proporcionada por los fondos de la UE ha permitido a las regiones más desfavorecidas invertir en infraestructuras, educación y fonnación, innovación y desarrollo empresarial, entre otros aspectos clave para el crecimiento económico. Además, la política regional h a fomentado la cooperación y el intercambio de buenas prácticas entre las regiones, lo que ha pennitido una mejor utilización de los recursos y un mayor aprendizaje entre regiones.

En el octavo Informe sobre la cohesión económica. social y territorial de la UE de la Comisión Europea, que se publica cada tres años se evalúa la convergencia real en la UE en términos de PIB per cápita empleo educación y pobreza, tanto entre Estados miembros como dentro de ellos. Entre sus conclusiones se puede destacar que, para que la política de cohesión pueda seguir desempeñando su función actual como vector de inversión y creación de empleo, como instrumento para reducir las disparidades regionales e intrarregionales y como mecanismo de solidaridad para todas las regiones de la Unión, debe contar con una financiación sólida basada en los principios de asociación y gobemanza multinivel.

Asimismo se concluye que la convergencia se ha visto impulsada por el fuerte creci miento de las regiones menos desarrolladas. pero es probable que los beneficios derivados de los menores costes y el rendimiento de sus inversiones disminuyan con el tiempo y que las regiones menos desarrolladas tendrán que fomentar la educación y la fonnación, aumentar la inversión en investigación e innovación y mejorar la calidad de sus instituciones, sin dejar de invertir en infraestructuras, para mantener un crecimiento constante y evitar caer en una trampa del desarrollo, acabar con la brecha de conectividad y garantizar el acceso a servicios de calidad y a unas condiciones de vida dignas.

Por otro lado, se propone el principio de "no perjudicar a la cohesión", que implica que ninguna acción debe obstaculizar el proceso de convergencia o aumentar las diferencias entre regiones sugiriendo la participación del Comité Europeo de las Regiones en el diseño de este principio.Finalmente, destaca que el PIB, único indicador del nivel de desarrollo, no tiene en cuenta lasostenibilidad medioambiental, la eficiencia en el uso de los recursos, la integración y el progresosocial recondando que, además de las cuestiones económicas, la salud, la educación, la sostenibilidad, la equidad y la inclusión social son parte integrante del modelo de desarrollo de la Unión. Para ello, plantea que el PIB se complemente con nuevos criterios (por ejemplo, de carácter social, medioambiental o demográfico) para ofrecer una mejor visión de conjunto de la situación socioeconómica de las regiones, abordar las actuales prioridades de la Unión, como el Pacto Verde Europeo y el pilar europeo de derechos sociales, y reflejar mejor las transiciones ecológica, digital y demográfica y el bienestar de las personas.

Para el estudio de la convergencia entre regiones, se pueden analizar varias variables que pueden indicar si se están reduciendo las disparidades económicas, sociales y territoriales entre las regiones. Estas variables son ampliamente reconocidas en la literatura económica y en los informes de la Comisión Europea sobre la política regional de la UE.

Para proporcionar más detalles y datos. específicos sobre la convergencia regional en la UE, se pueden utilizar fuentes como el Informe de Convergencia de la Comisión Europea la Base de Datos Regional de Eurostat y otros informes y estudios de organismos internacionales y de investigación económica.

Es importante destacar que cada análisis de convergencia regional puede utilizar diferentes fuentes y metodologías, y los resultados pueden variar según los datos y las variables utilizadas en el análisis. Algunas de estas variables son:

1. PIB per cápita: el PIB per cápita es una medida de la producción económica por persona en una región. Si las regiones más pobres están creciendo a un ritmo más rápido que las más ricas, es posible que haya una convergencia.
2. Tasa de empleo: la tasa de empleo es una medida de la proporción de la población en edad laboral que está empleada. Si las regiones más pobres están aumentando su tasa de empleo a un ritmo más rápido que las más ricas, es posible que. haya una convergencia en términos de mercado laboral.
3. Inversión en investigación y desarrollo ((+D): la inversión en l+D es una medida del gasto en investigación y desarrollo de una región. Si las regiones más pobres están aumentando su inversión en I+D a un ritmo más rápido que las más ricas, es posible que haya una convergencia en términos de innovación.
4. Nivel de educación: el nivel de educación es una medida del nivel de formación de la población en una región. Si las regiones más pobres están mejorando su nivel de educación a un ritmo más rápido que las más ricas, es posible que haya una convergencia en términos de capital humano.
5. Infraestructuras: la calidad de las infraestructuras en una región, como carreteras, ferrocarriles, aeropuertos y puertos, puede influir en el crecimiento económico y la competitividad de la región. Si las regiones más pobres están mejorando sus infraestructuras a un ritmo más rápido que las más ricas, es posible que haya una convergencia en términos de infraestructuras.

La convergencia real de la cohesión territorial y social dentro de la UE es un tema de investigación importante, que ha sido abordado por numerosos estudios. Estos estudios utilizan diferentes enfoques y metodologías para evaluar la convergencia real, tanto entre Estados miembros como dentro de los Estados miembros.

Los estudios empíricos sobre la convergencia real también utilizan diferentes enfoques econométricos, como el análisis de regresión, el análisis de series temporales y el análisis de panel. Estos estudios suelen evaluar la convergencia real en términos de PIB per cápita, ingresos, empleo, educación y otros indicadores de bienestar.

Por ejemplo, el estudio de \href{https://www.researchgate.net/publication/318285230_The_Impact_of_the_Regional_Policy_on_the_European_Economy}{LAGO-PEÑAS et al. (2019)} evaluó la convergencia real en términos de PIB per cápita y empleo en los Estados miembros de la UE durante el período 2000-2014, utilizando modelos de panel. Los autores encontraron una disminución en la disparidad del PIB per cápita entre los países estudiados durante el período estudiado, pero también encontraron una mayor divergencia en cuanto al empleo. Los autores utilizan una técnica econométrica conocida como "synthetic control" para comparar el desempeño económico de las regiones que han recibido financiación de la política regional de la UE con el de un grupo de regiones de control que no han recibido financiación.

Otro ejemplo es el estudio de Vassilis Monastiriotis (2021), que evaluó la convergencia real dentro de los Estados miembros de la UE en términos de ingresos, empleo y pobreza durante el período 2008-20 18, utilizando modelos de panel. El autor encontró que la convergencia real ha sido limitada en los Estados miembros de la UE y que ha habido una mayor divergencia en algunos indicadores de bienestar, (índice de Gini, nivel de ingresos, tasa de empleo. porcentaje de la población en riesgo de pobreza o tasa de pobreza relativa, entre otros) especialmente en países del sur de Europa.

Los estudios empíricos sobre la política de cohesión de la Unión Europea (UE) son numerosos y diversos en cuanto a sus enfoques y métodos de evaluación. A continuación, se presentan algunos de los hallazgos más relevantes de estos estudios. 

\href{https://onlinelibrary-wiley-com.sire.ub.edu/doi/10.1111/j.1468-2257.2011.00564.x}{LE GALLO (2011)} se centra en analizar los efectos de los Fondos Estructurales de la UE en el crecimiento económico regional. Específicamente, el autor examina si los Fondos Estructurales tienen un impacto más efectivo a nivel local o global en el crecimiento económico de las regiones.

El autor utiliza un modelo econométrico de panel para analizar los datos de crecimiento económico de las regiones europeas entre 1995 y 2004. El modelo incluye variables como el gasto de los Fondos Estructurales, el capital humano, la inversión privada y el tamaño de la población, entre otras. 

Los resultados del estudio muestran que el gasto de los Fondos Estructurales tiene un efecto positivo ·y significativo en el crecimiento económico de las regiones. Sin embargo, el autor encuentra que el impacto de los Fondos Estructurales es más efectivo a nivel local que a nivel global. Es decir, el gasto de los Fondos Estructurales tiene un mayor impacto en el crecimiento económico de las regiones más pobres y menos desarrolladas que en las regiones más ricas y desarrolladas.

El autor también examina el efecto del tamaño de la población en el impacto de los Fondos Estructurales en el crecimiento económico. Encuentra que las regiones más pequeñas se benefician más del gasto de los Fondos Estructurales que las regiones más grandes, lo que sugiere que los Fondos Estructurales tienen un mayor impacto a nivel local.

En general, el estudio sugiere que los Fondos Estructurales de la UE tienen un impacto positivo en el crecimiento económico regional, pero su efectividad depende del nivel de desarrollo de las regiones y del tamaño de su población. Además, el estudio destaca la importancia de considerar la dimensión local en la implementación de políticas regionales.

Por .otro lado, el estudio de \href{https://www.econstor.eu/handle/10419/120209}{Martín y Mayer (2008)}32 proporciona una evaluación empírica exhaustiva del impacto de la política regional de la UE en el crecimiento económico de las regiones europeas. Los autores examinan el impacto del gasto de los Fondos Estructurales de la UE en el crecimiento económico de las regiones utilizando diferentes especificaciones empíricas y medidas de crecimiento económico.

Aunque los autores encuentran que el gasto de los Fondos Estructurales tiene un efecto positivo y significativo en el crecimiento económico de las regiones menos desarrolladas, este efecto es relativamente pequeño.

En el estudio también destacan que el impacto del gasto de los Fondos Estructurales en el crecimiento económico de las regiones depende de la forma en que se mide el crecimiento económico. Por ejemplo, encuentran que el impacto es mayor cuando se utiliza la tasa de empleo como medida de crecimiento económico en lugar del PIB per cápita. Es decir, esto implica que el gasto de los Fondos Estructurales de la UE tiene un efecto positivo y significativo en el aumento de la tasa de empleo en las regiones menos desarrolladas de Europa. Por tanto, puede ser empleado como medida de crecimiento económico ya que proporciona información sobre la capacidad de una región para generar oportunidades laborales y reducir el desempleo, lo cual es crucial para el bienestar económico de la población y el desarrollo regional.

Este resultado puede deberse a que el gasto de los Fondos Estructurales está.dirigido a proyectos de infraestructura y desarrollo que pueden no tener un impacto inmediato en el PIB per cápita de las regiones, pero pueden generar empleo y mejorar las condiciones de vida de las personas y, por lo tanto, a mejorar su crecimiento económico. Por lo tanto, la tasa de empleo puede ser una medida más sensible en el corto y medio plazo para evaluar el impacto del gasto de los Fondos

Estructurales en el crecimiento económico regional. pero a largo plazo es posible que se observen efectos en el PIB per cápita a medida que las inversiones y el desarrollo impulsados por el gasto de los Fondos Estructurales se consoliden y se traduzcan en un crecimiento económico sostenible (pueden tener efectos positivos en la productividad y la capacidad de innovación de una región). En general, el estudio sugiere que la política regional de la UE ha tenido un impacto limitado en el crecimiento económico de las regiones y que se necesitan medidas adicionales para mejorar el impacto de la política regional de la UE en el desarrollo económico regional.

En otro estudio, "The employment impact of business investment incentives in declining areas: an evaluation of the EU 'Objective 2 Area' programs" de Andini et al. (2013), los autores utilizaron un enfoque econométrico de evaluación de impacto basado en el método de diferencia en diferencias. Esta metodología compara los cambios en el resultado de interés (en este caso, el empleo) en las regiones elegibles para los incentivos de inversión empresarial con las regiones no elegibles en el mismo período de tiempo. Los autores también controlaron por factores adicionales que podrían afectar el resultado, como las tendencias de empleo a nivel nacional y regional, las diferencias iniciales en el empleo y otras variables socioeconómicas. El análisis se centra en el programa Objetivo 2, que se dirigió a las regiones que no recibieron apoyo en el marco de la política de cohesión de la UE y encuentra que el efecto de los incentivos de inversión empresarial en las áreas en declive fue positivo pero modesto en términos de creación de empleo neto. Además, este efecto fue mayor en las regiones más ricas que en las más pobres y dependía del tipo de industria y del tamaí'lo de la empresa. Las empresas más grandes y las industrias con mayor productividad y salarios más altos fueron las que tuvieron el mayor impacto en la creación de empleo.

A continuación se resume en una tabla algunas de las conclusiones más relevantes en el análisis de la convergencia real de la política de cohesión territorial, social y económica de la Unión Europea tanto entre Estados miembros como dentro de los mismos:

\begin{center}
\begin{tabular}{|p{0.25\linewidth}|p{0.65\linewidth}|}
\hline
\textbf{ESTUDIO} & \textbf{CONCLUSIONES} \\
\hline
LAGOS-PEÑAS et al. (2019) & Hubo convergencia real entre los Estados Miembros de la UE en términos de PIB per cápita. Sin embargo, en cuanto a la convergencia real dentro de los Estados Miembros, es decir, la reducción de las disparidades económicas entre las regiones o áreas dentro de un mismo país, el estudio no encontró una convergencia significativa. Esto implica que las diferencias económicas entre las regiones dentro de los Estados Miembros se mantuvieron o incluso se ampliaron. \\
PINTERA (2021) & Se aplicó una prueba de convergencia desarrollada por PHILIPSS y SUL (2009) y no se encontró una convergencia general entre las regiones de la Unión Europea después de la crisis de 2008, pero se identificaron cinco clubes de convergencia en tasas de crecimiento de ingresos. \par Se observaron significativas disparidades en la membresía de los clubes dentro de los países (que convergen en tasas de crecimiento de ingresos en lugar de nivel de ingresos), indicando alta desigualdad económica y una división centro-periferia en la UE (visible en el contraste entre la capital y la periferia, ya que las capitales suelen converger hacia los clubes más altos tanto en los países antiguos como en los nuevos de la UE). \\
\hline
MARTIN y MAYER (2008) & La política de cohesión de la UE ha tenido un impacto positivo pero modesto en el crecimiento económico, especialmente en las regiones menos desarrolladas. \\
\hline
CHECCHI y GARCÍA-PEÑALOSA (2020) & Hubo convergencia real en términos de PIB per cápita y productividad en la UE entre 1995 y 2017, pero la desigualdad de ingresos se mantuvo constanteo aumentó en algunos países. \\
\hline
FRATESI y RODRÍGUEZ-POSE (2018) & La política de cohesión de la UE ha tenido un impacto positivo pero limitado en la convergencia real en términos de PIB per cápita y empleo en las regiones menos desarrolladas. Se necesitan políticas más específicas con un enfoque más adaptado y personalizado que se ajusten a las necesidades y oportunidades de cada región, con el objetivo de promover un desarrollo económico más equilibrado y una mayor convergencia regional.\\
\hline
BACHTLER y MARTINS (2010) & La política de cohesión de la UE ha contribuido a la convergencia real en términs de PIB per ápita y empleo en las regiones menos desarrolladas, pero la coordinación insuficiente y la falta de capacidad de gestión han sido obstáculos para una implementación efectiva. \\
ANDINI et al. (2013) & Los incentivos de inversión empresarial en las áreas en declive del programa \textit{Objetivo 2} de la UE tuvieron un impacto positivo en el empleo a corto plazo, pero no tuvieron un impacto significativo en el largo plazo. Se necesitan políticas más integrales para abordar los desafíos estructurales de las regiones en declive. \\
LE GALLO (2011) & La política de cohesión de la UE ha tenido un impacto positivo pero limitado en la convergencia real en términos de PIB per cápita y empleo en las regiones menos desarrolladas. La inversión en infraestructura y capital humano es importante para lograr una cobergencia real efectiva.\\
\hline
\end{tabular}
\end{center}

Es importante terter en cuenta que los estudios pueden utilizar diferentes metodologías y periodos de tiempo, lo que puede afectar a sus conclusiones. Sin embargo, en general, parecen coincidir en que la convergencia económica ha sido limitada en los últimos años, ·tanto dentro de los estados miembros como entre ellos. 

Por último, cabe hacer referencia a dos dificultades principales a las que se enfrenta cualquier intento de evaluar el impacto de la política regional:\footnote{%
    Bachtler. J., & Mendez.. C. (2019). EU Cohcsion Policy and the Role of the Regions: lnvestigating Synergies, Strategic Orientations and the Added Value of EU lnterventions. European Planning Studies, 27(2), 1 9 1 -204. Rodríguez-Pose, A., & Tselios, V. (2009). Mapping regional personal income distribution in Western Europe . .loumal of Economic Geography, 9(6), 773-797. Crescenzi, R., & Rodríguez-Pose. A. (2012). An •integrated' framework for the comparative analysis ofthe ten-itorial innovation dynamics of developed and emerging countries. Joumal of Economic Surveys. 26(3 ). 5 17-533. McCann. P .. & Onega-Argilés, R. (2015). Smart specialisation, regional growth and applications to European Union cohesion policy. Regional Studics. 49(8). 1291- 1302.
}

La primera dificultad es la falta de datos precisos y comparables en el tiempo y entre las regiones. Existen diferencias importantes en los sistemas estadísticos y de información en los países de la UE, lo que hace dificil comparar los datos de diferentes regiones. Además, los indicadores utilizados para medir la convergencia real no son siempre los mismos, lo que dificulta la comparación de los resultados de diferentes estudios.

La segunda dificultad es la identificación del efecto causal de la política regional en el crecimiento y la convergencia. Existen múltiples factores que influyen en el crecimiento económico y la convergencia, como la innovación, la educación, la infraestructura, la inversión privada, etc. Por lo tanto, es dificil aislar el efecto específico de la política regional en el crecimiento y la convergencia. Además, el impacto de la política regional puede ser afectado por factores externos, como las fluctuaciones del ciclo económico o los cambios en el contexto político y económico de la UE y sus estados miembros.

\clearpage
\section{Unión Europea: Política Social y de Empleo}
\puntosclave{
    Uno de los principales objetivos de la UE es promover el progreso social y mejorar las condiciones de vida y de trabajo de toda la ciudadanía de la UE. un principio consagrado en el preámbulo del Tratado de Funcionamiento de la Unión Europea (TFUE) y desarrollado a lo largo de diferentes apartados del Tratado.
    
    La UE sólo tiene competencias limitadas en el ámbito social, por lo que la responsabilidad del empleo y las políticas sociales recae principalmente en los gobiernos nacionales.
    
    La UE toma medidas para garantizar la coordinación de las políticas de empleo de los Estados miembros, en particular definiendo las orientaciones de dichas políticas y alentando a los Estados miembros a compartir las mejores prácticas en sectores como la inclusión social, la pobreza y las pensiones.
    
    Una de las bases de la política social de la UE es el Pilar europeo de derechos sociales.
    
    El Fondo Social Europeo Plus (FSE+), es el principal instrumento de la UE para apoyar la aplicación del Pilar europeo de los derechos sociales y alcanzar los tres objetivos principales de la UE para 2030.
}

Hemos presentado y analizado la importancia de la Política de Cohesión, cómo esta se ha construido a lo largo de los años, su faceta multidimensional actual y los resultados que, aparentemente, ha tenido. Sin embargo, una visión completa de la proyección social del proyecto europeo exige que presentemos otra de las grandes políticas comunitarias en la materia: la Política Social y de Empleo. 

La Política Social y de Empleo encuentra su fundamento jurídico en el artículo 3 del TUE y los arts. 9, 10, 19, 45 a 48 y 145 a 161 del TFUE. Constituye, por tanto, una pilar fundamental del proyecto de integración.

\textbf{¿Qué queremos decir?} A lo largo de esta sección, presentaremos....

Para ser claros, existen dos formas de abordar esta política comunitario. Por un lado estaría su vertiente jurídico, cuyo objetivo esencial es buscar la efectiva materialización de la \textbf{ciudadanía europea}. En este, trataríamos las sucesivas barreras a la libre circulación de personas, a la libre residencia, el Espacio de Libertad, Seguridad y Justicia de la Unión Europea, la jurisprudencia del TJUE relatica a la Carta Europea de Derechos Humanos, etc. No trataremos esta vertiente aquí puesto que, aún siendo la más extendida y una de las grandes virtudes (sino la que más) del proyecto europeo, o bien ya se ha tratado en otros tema [\authorfont{Ver Tema 3.B.}] o bien excede del objeto del tema y la oposición, puesto que sería una aproximación más jurídica que económica. 

Por otro lado, que sí es objeto del tema y pasamos a tratar a continuación, tendríamos las políticas comunitarias llevadas a cabo sobre la cohesión social y el mercado laboral o la Seguridad Social. En este sentido, el artículo 3 del TUE establece que la Unión \textbf{tiene el deber de aspirar al pleno empleo y trabajar por el progreso social}. El fomento del empleo, la mejora de las condiciones de vida y de trabajo, la protección social, el diálogo entre la dirección y el resto del personal, el desarrollo de los recursos humanos para garantizar un nivel de empleo elevado y duradero y la lucha contra la exclusión social son los objetivos comunes de la Unión y sus Estados miembros en materia social y laboral, tal y como se establece en el artículo 151 del TFUE.

Su fundamento jurídico se encuentra recogido tanto en el TUE como el TFUE. En el TFUE, se establece que, en todas las políticas y acciones, la Unión tendrá en cuenta las exigencias relacionadas con la promoción de un nivel de empleo elevado, con la garantía de una protección social adecuada, con la lucha contra la exclusión social y con un nivel elevado de educación, formación y protección de la salud humana. Asimismo, se enumera los objetivos de la política social de la Unión:\footnote{Art. 9 TFUE y art. 151 TFUE.}
• El fomento del empleo.
• La mejora de las condiciones de vida y de trabajo.
• Una protección social adecuada.
• El diálogo social.
• El desarrollo de los recursos humanos para garantizar un nivel de empleo elevado y duradero.
• La lucha contra la exclusión.

\subsection{Desarrollo histórico: ¿El Modelo Social Europeo?}
El denominado modelo social europeo, a diferencia del norteamericano o del japonés, ha sido descrito en la literatura comparada como un desarrollado sistema de redes de protección social, que ha configurado, en última instancia, la apuesta europea por el Estado del Bienestar (VIÑALS, 2005). Ahora bien, a pesar de ser masiva su presencia en los países europeos, en él subyace una gran heterogeneidad y multiplicidad de valores, contenidos, esquemas institucionales, rendimientos y modelos que funcionan como espacios \textit{con trazos básicos comunes en las políticas sociales4 de sus Estados}, aunque con grandes \textit{dosis de pluralidad interna} (ADELANTADO y GOMÁ, 2000). Precisamente, esta ausencia de uniformidad en los Estados del Bienestar, unida al gran protagonismo alcanzado por las políticas sociales nacionales y al hecho de que la Comunidad Europea naciera como un proceso de integración económica, han condicionado, en buena medida, el desarrollo de una plena y activa Política Social Comunitaria.\footnote{%
    Su evolución no ha estado exenta de incertidumbres y dificultades, vinculadas en la mayor parte de las ocasiones a la carencia de un marco institucional adecuado (BAILEY, 2008), habiendo sido, históricamente, fuente de recelos permanentes entre los países miembros, por considerar a la política social como una expresión de su soberanía nacional.
}

No obstante, este proceso de formación, amplio e irregular, hace que se pueda hablar de distintas etapas en el desarrollo histórico de la política social de la Unión Europea, previas a la fase por la que atraviesa en la actualidad y que será analizada más adelante. 

\subsubsection{Concienciación comunitaria: 1958-1973}
Dado que, de las dos posturas que subyacen en el debate teórico reinante sobre la política social (convergencia social espontánea ex post versus homogenización ex ante), el Tratado de Roma eligió la de aquéllos que consideran que el progreso social se producirá de forma espontánea y automática, con el libre funcionamiento del mercado y la eliminación de los obstáculos a la libre competencia, no resulta extraño que algunos autores (MOSLEY, 1990) hayan denominado a los años de 1958-1973 el periodo del \textit{olvido} y otros (MUÑOZ DE BUSTILLO y BONETE, 2009), el de la \textit{concienciación} comunitaria en materia social.\footnote{%
    Los orígenes de la Política Social y de Empleo de la Unión Europea se remontan a la constitución de la CECA (1952). Su funcionamiento preveía que los trabajadores se beneficiaran de la \textbf{ayuda de readaptación}. Esta ayuda se concedía a los trabajadores de los sectores del carbón y del acero cuyos puestos de trabajo se veían amenazados por la reestructuración industrial. Más tarde, con la creación de la Comunidad Económica Europea mediante el Tratado de Roma (1957) y el impulso integrador que este generó, la Política Social y de Empleo adquirió una profundidad mucho mayor. Por un lado, preveía la coordinación de los Sistemas de Seguridad Social de los Estados miembros, para permitir a los trabajadores y a sus familias aprovechar al máximo el derecho a la movilidad y a buscar un empleo libremente en todo el mercado común. Tambiém consagró el principio de igualdad de retribución para hombres y mujeres, que el Tribunal de Justicia de la Unión Europea reconoció como directamente aplicable. Además, mediante la constitución del Fondo Social Europeo (FSE) se buscada dotar a la Comisión de fondos necesarios para luchas contra el desempleo de manera coordinada.
} ¿Por qué? Durante este periodo los seis países fundadores atravesaran una situación económica de altas tasas de crecimiento y bajas tasas de desempleo, la relativa homogeneidad de sus políticas sociales nacionales, junto a las estrechas, aunque convergentes, diferencias de costes sociales existentes entre ellos, contribuyeron a tal elección.

La Cumbre de París de 1972 sentaría, no obstante, las bases para que la política social pasara a ser considerada como un elemento positivo en la transformación comunitaria, al exigir los Jefes de Estado y de Gobierno allí reunidos un enérgico avance en su ámbito de actuación, definir un conjunto de objetivos prioritarios a lograr y encargar a la Comisión la elaboración de un programa social.

\subsubsection{Hacia una Política Social activa (Espacio Social Europeo): 1974-1987}
Fruto de este mandato fue la adopción, en 1974, del primer Programa de Acción Social (1974-1976) que reconocía que la política social comunitaria debía \textit{cumplir  una función propia y proporcionar una contribución esencial para la realización de sus objetivos}. Expresaba, asimismo, la voluntad política de adoptar las medidas necesarias para alcanzar los objetivos siguientes: pleno empleo y mejor; la mejora de las condiciones de vida y de trabajo; y la participación de los empresarios y trabajadores en las decisiones socioeconómicas de la Comunidad y de estos últimos en la vida de las empresas.

Así, entre 1975 y 1980, se aprobaron directivas relacionadas con la aplicación del principio de igualdad de trato entre hombres y mujeres, la coordinación de los regímenes nacionales de seguridad social para las personas que se desplazan dentro de la Comunidad, o con la salvaguardia de los derechos de los trabajadores ante despidos colectivos, insolvencia del empresario y traspasos de empresas o centros de actividad; se adoptó un programa específico orientado hacia la pobreza; se mejoró la coordinación de los servicios nacionales de empleo; y se hicieron propuestas para incorporar a los jóvenes, discapacitados y desempleados a la vida profesional. La progresiva consolidación del proceso de integración europea y la creciente heterogeneidad de los regímenes de bienestar de los Estados miembros actuaron de catalizadores de este modesto, aunque incipiente, giro social de los 70.

Sin embargo, la persistencia de graves dificultades económicas, la influencia de las doctrinas conservadoras de neutralización del intervencionismo social, lideradas por el Reino Unido, y el avance de las posiciones euroescépticas condujeron a una relativa parálisis legislativa en la política social de los años 1981-8721 (ADELANTADO y GOMÁ, 2000), aunque no lograron impedir que la dimensión social fuera ganando importancia en Europa bajo la influencia de la Comisión DELORS.

De esta forma, la propuesta del gobierno francés que emana de las elecciones de 1981, sobre la necesidad de un Espacio Social Europeo para potenciar el papel de la política social comunitaria en la construcción del Espacio Económico Europeo, tendría su reflejo en un nuevo Programa de Acción Social Comunitario (1984), en el Libro Blanco de la Comisión sobre la culminación del mercado interior (1985) y en el Acta Única Europea (1986).\footnote{%
    El primero contemplaba acciones e iniciativas en los ámbitos del empleo, los aspectos sociales de las nuevas tecnologías, la formación, la protección social, la evolución demográfica y el diálogo social europeo, para asegurar la dimensión social del gran mercado interior que habría de alcanzarse en 1992. El segundo daba cuenta de la interacción existente entre el mercado interior y la política social, así como de las oportunas medidas que aquél requería para alcanzar los objetivos en materia de empleo y de seguridad social. El tercero introducía la trascendental noción de la cohesión económica y social de la Comunidad.
}

\subsubsection{Cohesión económica y \textit{geometría variable}: 1988-1992}
1988 y 1989 serían años importantes para la política social europea, aunque sus logros y avances más bien modestos. La reforma de los Fondos Estructurales de 1988, intensa y global, derivada del AUE, permitió al FSE: (i) financiar en exclusiva las acciones elegibles de los nuevos objetivos 3 (luchar contra el desempleo de larga duración) y 4 (facilitar la inserción profesional de los jóvenes); (ii) intervenir (junto al FEDER y FEOGA-Orientación) en las actuaciones de los objetivos 1 (fomentar el desarrollo y el ajuste estructural de las regiones menos desarrolladas), 2 (reconvertir las regiones gravemente afectadas por el declive industrial) y 5b (fomentar el desarrollo de las zonas rurales), con una acentuada dimensión social a partir de entonces; y, (iii) co-subvencionar algunas de las numerosas iniciativas impulsadas por la Comisión Europea para dar solución a problemas específicos de dimensión europea.

Por otra parte, la formulación, en 1989, del \textit{Programa de Acción Social} (1989-1994), basándose en la Carta Comunitaria de los Derechos Sociales Fundamentales de los Trabajadores, contempló $47$ medidas sobre el derecho a la libre circulación, el empleo y la formación profesional, la regulación de las condiciones de trabajo, la protección social y los derechos colectivos laborales. Sin embargo, la negativa actitud del Reino Unido (y de la patronal europea) ante los compromisos recogidos en la Carta; los importantes cambios políticos que acontecen en el centro-este de Europa (caída del muro de Berlín y de los regímenes comunistas circundantes); las consecuencias económicas de la reunificación de Alemania; el estrecho margen de maniobra de las perspectivas financieras (1988- 1992) de la Comunidad; y las difíciles negociaciones del Tratado de Maastricht limitaron, en buena medida, su aplicación (BRUNET, 1999).\footnote{%
    Entre las disposiciones aprobadas figuran: las directivas sobre salud y seguridad en el lugar de trabajo (1989), protección de los trabajadores contra los riesgos derivados de la exposición a agentes biológicos (1990), seguridad y salud en el trabajo de los empleados con contratos temporales (1991), disposiciones mínimas sobre señalización para la seguridad y salud en el trabajo (1992) y protección de la mujer embarazada, que haya dado a luz o esté en periodo de lactancia que trabaja (1992); los programas de acción para el desarrollo de una formación profesional continuada (1990) y para la igualdad de oportunidades entre mujeres y hombres (1991); las recomendaciones de la Comisión sobre la creación de criterios comunes para la suficiencia de renta y la extensión de la asistencia social en los sistemas de protección social (1992) y sobre la convergencia en las políticas y los objetivos de protección social (1992), etc.
}

\subsubsection{La Europa Social y del Empleo: 1993-1999}
Las debilidades que presentaba la economía europea en 1993, la extraordinaria magnitud que había alcanzado el nivel de desempleo en su seno y el horizonte de la UEM ralentizaron la voluntad política de completar el gran mercado interior con una fuerte dimensión social y forzaron una reflexión acerca de cuáles deberían ser las políticas a adoptar para alcanzar un desarrollo que permitiera hacer frente a la competencia internacional (LÓPEZ, 2009).\footnote{%
    Tres documentos publicados por la Comisión sobresalen en tal sentido:
\begin{lnum}
    \item El Libro Verde sobre el futuro de la política social europea (Comisión de las Comunidades Europeas, 1993a). Dio lugar a un extenso proceso de consulta y concluyó reconociendo que el desempleo en Europa se había convertido en un problema estructural y que debía tratarse como tal;
    \item El Libro Blanco sobre crecimiento, competitividad y empleo (Comisión de las Comunidades Europeas, 1993b). Presentaba, básicamente, las grandes orientaciones económicas: la consecución de una economía sana, abierta, descentralizada, más competitiva y solidaria y la realización de cambios profundos en la política de empleo; y,
    \item El Libro Blanco sobre la política social (Comisión de las Comunidades Europeas, 1994). Hacía hincapié en la cuestión de que la competitividad y el progreso social podían ir unidos, a la vez que fijaba sus grandes objetivos: el empleo, la necesidad de una sociedad que exija la colaboración de todos y las modalidades de desarrollo del fundamento jurídico.
\end{lnum}
}

En este contexto surge la denominada \textbf{Estrategia Europea de Empleo}, que alcanzaría su máxima representatividad en el Tratado de Ámsterdam.\footnote{%
    Se adoptó tras ser debatida en los sucesivos Consejos Europeos de Essen (1994), Madrid (1995), Florencia (1996) y Ámsterdam (1997). El Tratado de Amsterdam finalmente incorporó el nuevo Título dedicado al empleo; asumir competencias comunitarias en esta materia; constitucionalizar los derechos sociales de 1961 y 1989 y habilitar la política social a todos los Estados miembros.
} De hecho, los mandatarios europeos decidieron no esperar hasta la entrada en vigor del Tratado para lanzar la estrategia y en el Consejo Extraordinario sobre el Empleo celebrado en Luxemburgo, en noviembre de 1997, acordaron poner en marcha, en 1998, una serie de orientaciones basadas en análisis de situación y sustentadas en cuatro pilares fundamentales: (i) el desarrollo del espíritu empresarial; (ii) la mejora de la capacidad de inserción profesional (empleabilidad); (iii) la reactivación de la capacidad de adaptación de empresarios y trabajadores al cambio económico y tecnológico (adaptabilidad); y, (iv) la potenciación de la igualdad de oportunidades entre mujeres y hombres y para las personas discapacitadas.\footnote{%
    Adoptaron, igualmente, un programa de trabajo continuo de planificación, vigilancia, examen y reajuste anual de las políticas nacionales de empleo (conocido como el Proceso de Luxemburgo), que se vería reforzado, posteriormente, por el denominado Pacto para el Empleo en el Consejo Europeo de Colonia de 1999.
}

El \textit{Programa de Acción Social a Medio Plazo 1995-1997}, para hacer frente a las necesidades y desafíos planteados, orientó sus propuestas hacia el empleo (su primera prioridad); la educación y la formación (factores de suma importancia para la estabilidad social comunitaria); la instauración de un mercado europeo del empleo; la promoción de mejores normas de empleo; la igualdad de oportunidades entre mujeres y hombres; la política y promoción sociales (buscando una sociedad activa para todos); la salud pública; el desarrollo de la dimensión internacional; la reactivación de los agentes sociales; el análisis e investigación de la política social a medio y largo plazo; y la aplicación más eficaz del derecho europeo. Todas estas medidas financiadas mediante el Fondo Social Europeo. Más adelante, el Programa de Acción Social 1998-2000, apoyándose en los resultados del anterior e intentando aprovechar el impulso social de la Estrategia de Empleo y del Tratado de Ámsterdam, redujo sus prioridades a las tres siguientes: (i) empleo, capacitación y movilidad (creación de empleo, supresión del paro y fomento de la libertad de circulación de los trabajadores); (ii) cambios en el mundo del trabajo (modernización de la organización laboral y fomento de la adaptabilidad, crear un lugar de trabajo sano y seguro, adelantarse al cambio industrial y aprovechar las oportunidades que ofrece la sociedad de la información); y, (iii) desarrollo de una sociedad no excluyente (fomentar la inclusión social y la creación de una sociedad sana, modernizar y mejorar la protección social, y conseguir la igualdad y combatir la discriminación).


\subsection{Estructura actual: Pilar Europeo de Derechos Sociales (2017)}
Como vemos, los asuntos en materia de empleo y política social en la Unión Europea han ido adquiriendo mayor relevancia a medida que se ha ido profundizando el proceso de integración europea. Además, estos ámbitos están, en muchos casos, ligados a la Política de Cohesión presentada en el apartado anterior. Sin embargo, por su propia naturaleza podemos haber apreciado que consiste en un conjunto de actividades mucho más heterogéneas que la Política anterior. Esto se debe a que, aunque parezca una linea de Política única, en la realidad consiste en dos lineas de acción diferentes, con grades diferencias de aproximación y consenso. 

De hecho, esto queda claramente diferenciado, además de por la evolución histórica presentada, cuando entramos a analizar la normativa en la materia:\footnote{%
    Además, en relación con ambos ámbitos cada recordar que el TFUE establece que \textit{en la definición y ejecución de sus políticas y acciones, la Unión tendrá en cuenta las exigencias relacionadas con la promoción de un nivel de empleo elevado, con la garantía de una protección social adecuada, con la lucha contra la exclusión social y con un nivel elevado de educación, formación y protección de la salud humana.}
}
\begin{la}
    \item \textbf{Política Social}. Por un lado, se puede afirmar que la Política Social tiene una categoría propia en la delimitación de competencias: no es ni enteramente compartida, ni las funciones de la UE son de apoyo, coordinación o complemento.\footnote{%
        De hecho, en el TFUE, el artículo sobre la capacidad de la UE para tomar iniciativas en política social se encuentra después del artículo sobre competencias compartidas e inmediatamente antes del de ámbitos en los que la UE sólo tiene competencias de apoyo, coordinación y complemento. Sin embargo, el articulado del Tratado no llega a explicitar dichos aspectos, de manera que el conjunto de ámbitos objeto de competencia compartida depende de la interpretación que se hace del articulado y la jurisprudencia. Un ámbito que se suele poner como ejemplo de competencia compartida es el de igualdad de retribución entre trabajadores.
    } De acuerdo con algunos autores, la explicación de esta particular categorización de la competencia en Política Social es de carácter político;
    \item \textbf{Política de Empleo}. Por otro lado, y en cuanto a la Política de Empleo, el TFUE establece que la Unión tomará \textbf{medidas de coordinación}, en particular \textbf{definiendo las orientaciones de dichas políticas}. En este caso, aunque no es una competencia compartida ni complementaria, se podría considerar que el papel asignado a la UE es algo más relevante que la Política Social, pues aquí tomará medidas para garantizar la coordinación y definirá orientaciones, mientras que en la política social simplemente podrá tomar iniciativas. 
\end{la}

La razón de este comentario es dejar claro que resulta muy, muy, dificil presentar su estructura (tanto actual como histórica, como se abrá podido obsevar) de una manera coherente, equiparable a la anterior. La razón es que la actividad comunitaria en ambas políticas no solo difiere mucho en sus medios materiales, sino también en sus objetivos, funciones y dimensiones. Para entenderlo de manera sencilla, podemos decir que no existe como tal una estrategia integrada y coherente que pueda servir para establecer paralelismos entre las acciones realizadas en una y en otra. 

No obstante, con el objetivo de presentarlo de manera coherente, trazaremos un paralelismo con la Político de Cohesión e intentaremos identificar los principales dimensiones que configuran la Política Social y de Empleo de la UE basándonos en aquellas que juegan un papel fundamental en la ejecución del Pilar Europeo de Derechos Sociales, adoptado por la Comisión en 2017; y cuyo Plan de Acción se presentó en 2021.

Este Pilar presenta veinte objetivos estratégicos, articulados como principios, que pueden divirse en tres categorías:\footnote{%
    \textcolor{red}{OJO ->} Se puede aclarar al Tribunal que en ningún momento vamos a hacer mención de los instrumentos jurídicos comunitarios que se han elaborado y adoptado a lo largo de los años en estas materias concretas. Si bien es cierto que no debería suponer un problema, la realidad es que una grandísima parte de la Política de Empleo (y, en menor medida, de la Política Social) se materializa mediante esta dimensión jurídica. Si se cree necesario, se puede decir durante la exposición o, incluso, en las conclusiones. Para ver la normativa sobre esta materia: [\authorfont{Ver Anexo \ref{anx:normativa_empleo}}]
}
\begin{la}
    \item \textbf{Capítulo I: Igualdad de oportunidades y de acceso al mercado de trabajo}. Incluye los aspectos de la equidad relacionadas con la educación, las capacidades y el aprendizaje permanente, el género, las desigualdades y la movilidad social, las condiciones de vida y la juventud;
    \item \textbf{Capítulo II: Mercados de trabajo dinámicos y condiciones de trabajo justas}. Responde a la pregunta de si los mercados funcionan de forma eficiente y libre para ayudar a la reubicación y la búsqueda de empleo en una economía competitiva. Los indicadores a este respeto miden el funcionamiento de los mercados de trabajo, el apoyo al empleo y las transiciones y, por otra parte, la dimensión de equidad en las condiciones salariales y de trabajo; y,
    \item \textbf{Capítulo III: Apoyo público/protección e inclusión social}. Tiene por objeto los resultados de equidad obtenidos por la acción pública o la protección social. Este conjunto de indicadores se refiere principalmente a la prestación de servicios y a las redes de seguridad social.
\end{la}  

\begin{specialblock}{Pilar Europeo de Derechos Sociales: Situación actua (2026)}

    \textcolor{red}{OJO ->} Este Bloque está pensado para introducir las últimas novedades relativas a la aplicación e implementación del Plan de Acción del Pilar Europeo de Derechos Sociales. Podemos consultar la evolución histórica de su implementación en: [\authorfont{Ver Anexo \ref{anx:consolidacion_peds}}]

    
    
\end{specialblock}


\subsubsection{Dimensión financiera}
La primera de la que deberíamos hablar sería la dimensión financiera. Al igual que en el caso anterior, existen diferentes programas, o Fondos, de financiación de esta Política comunitaria.

Sin embargo, y a diferencia del anterior, la clasificación no sigue una linea estructural y/o finalista, sino que más bien se distinguen por su ambición, cuantía y marco temporal de sus objetivos.

\paragraph{Fondo Social Europeo Plus (FSE+, )}
Sin duda el más importante es el Fondo Social Europeo (FSE), que pasa a llamarse Fondo Social Europeo Plus (FSE+) en 2017. Tiene por objeto lograr elevados niveles de empleo y una protección social justa, y formar a una mano de obra cualificada y resiliente, preparada para la transición a una economía ecológica y digital. También tiene por objeto promover la movilidad laboral, la educación y la inclusión social, y construir sociedades inclusivas y cohesionadas capaces de erradicar la pobreza. 

Es, sin duda, el principal instrumento de la Unión dedicado a invertir en las personas, reuniendo un conjunto de programas y Fondos de muy diversa índole.\footnote{%
    Por ejemplo, la Iniciativa de Empleo Juvenil, el \href{https://employment-social-affairs.ec.europa.eu/policies-and-activities/funding/fund-european-aid-most-deprived-fead_es}{Fondo de Ayuda Europea para las Personas Más Desfavorecidas} (FEAD) y el \href{https://european-social-fund-plus.ec.europa.eu/en/direct-management-easi}{Programa de la Unión Europea para el Empleo y la Innovación Social} (EaS).
} Los recursos financieros se canalizan a través de los Estados miembros y las regiones, por lo que no financia proyectos de forma directa. La Comisión Europea acuerda la cantidad por objetivo y país (las regiones objetivo se establecen en el marco de la Política Regional de la UE) y aprueba las prioridades de la programación de cada Estado miembro. En lineas generales, los programas suelen seguir concentraciones temáticas cubiertas en las Recomendaciones Específicas por País (REP) del Semestre Europeo (que luego veremos), así como en sus Programas Nacionales de Reforma.

El Fondo Social Europeo Plus es la continuación del Fondo Social Europeo, que hemos visto en el apartado anterior. Su cambio se debe, fundamentalmente a la incorporación de una visión más holística de las estrategias de cohesión, incorporando también los retos relativos a aspectos sociales, como la marginalidad, la pobreza o la exclusión, y de empleo, como la formación o el trabajo juvenil, en linea con los \href{https://commission.europa.eu/document/download/e03c60e7-4139-430b-9216-3340f7c73c20_es?filename=social-summit-european-pillar-social-rights-booklet_es.pdf}{Principios del Pilar Europeo de Derechos Sociales, adoptado en 2017}.

En 2021 se adoptó también la Recomendación del Consejo sobre el establecimiento de una Garantía Infantil Europea para prevenir y combatir la exclusión social garantizando el acceso de los niños y niñas necesitados a la educación y los cuidados de la primera infancia, la educación, la asistencia sanitaria, la nutrición y la vivienda. En el marco del Fondo Social Europeo Plus (FSE+) 2021-2027, todos los Estados miembros deben asignar un porcentaje adecuado de sus recursos para combatir la pobreza infantil. En el caso de los países con una tasa de niños y niñas en riesgo de pobreza o exclusión social superior a la media de la Unión, esta asignación debe ascender al $5\%$ del FSE+. Para respaldar la aplicación de la Garantía, la Comisión creó un nuevo marco de seguimiento en enero de 2024, que se actualizó de nuevo en mayo de 2025 y que debía aplicarse por parte de los Estados miembros y durante el Semestre Europeo. En mayo de 2026, la Comisión pondrá en marcha una nueva iniciativa para reforzar la Garantía Infantil, con motivo de su quinto aniversario.

\paragraph{Fondo Europeo de Adaptación a la Globalización (FEAG, 2007)}
El Fondo Europeo de Adaptación a la Globalización (FEAG) se creó en el MFP 2007-2013 para apoyar al personal despedido debido a los cambios estructurales en el comercio mundial derivados de la globalización. Mientras que el FSE+ apoya programas dirigidos a alcanzar los objetivos a largo plazo de mantener o reinsertar a las personas en el mercado laboral, el FEAG responde a emergencias específicas, como despidos masivos, durante períodos de tiempo limitados.

El Fondo Europeo de Adaptación a la Globalización para Trabajadores Despedidos (FEAG) es un instrumento especial de la Unión Europea para expresar la solidaridad de la misma con las personas europeas que trabajan por cuenta ajena o por cuenta propia que han perdido su puesto de trabajo como consecuencia de reestructuraciones, y ayudarles a encontrar un nuevo empleo.

El Fondo Europeo de Adaptación a la Globalización para Trabajadores Despedidos (FEAG) contribuye a la creación de una economía europea más dinámica y competitiva mejorando las capacidades y la empleabilidad de las personas que se han quedado sin empleo, a fin de ayudarles a encontrar un mejor trabajo.

Como norma general, el Fondo Europeo de Adaptación a la Globalización para Trabajadores Despedidos (FEAG) puede activarse cuando una única empresa (incluidos sus proveedores y los transformadores de sus productos) despide a más de 200 personas y también puede ser solicitada por pymes de diversos sectores de la misma región o de un sector concreto en una o más regiones vecinas.

Este fondo tiene un presupuesto anual de 210 millones de euros para el período 2021-2027. Puede financiar entre el 60 % y el 85 % del coste de los proyectos concebidos para ayudar a los trabajadores despedidos a encontrar un nuevo empleo o a crear su propia empresa.

Las intervenciones del Fondo Europeo de Adaptación a la Globalización para Trabajadores Despedidos (FEAG) las llevan a cabo las administraciones nacionales o regionales. Cada proyecto tiene una duración de dos años.


\subsubsection{Dimensión Política}
En base a la adopción del Pilar Europeo de Derechos Sociales, y con el fin de materializar lo que en el se expone, el 4 de marzo de 2021 la Comisión Europea publicó su \href{https://op.europa.eu/webpub/empl/european-pillar-of-social-rights/es/}{Plan de Acción del Pilar Europeo de Derechos Sociales}, una guía para la implementación de los principios del Pilar a todos los niveles: europeo, nacional, regional y local, porque según la propia Comunicación de la Comisión: Los veinte principios del pilar europeo de derechos sociales son el faro que nos guía hacia una Europa social fuerte y que determinan el enfoque del nuevo \textit{código normativo social}.\footnote{%
    Este Plan de Acción se ha basado en una consulta a gran escala iniciada hace casi un año en la que participaron, con sus contribuciones, ciudadanos, instituciones y organismos de la UE, Estados miembros, autoridades regionales y locales, interlocutores sociales y organizaciones de la sociedad civil.
} En este Plan de Acción se establece una serie de acciones de la Unión que la Comisión se compromete a adoptar, y se propone tres objetivos a escala de la Unión que deben alcanzarse para 2030 y que ayudarán a dirigir las políticas y reformas nacionales:
\begin{la}
    \item Al menos el $78\%$ de las personas de entre $20$ y $64$ años debe tener trabajo;
    \item Al menos el $60\%$ de los adultos debe participar en actividades de formación cada año;
    \item El número de personas en riesgo de pobreza o exclusión social debe reducirse en al menos $15$ millones, y de esos al menos $5$ millones deberían ser menores.
\end{la}

Hacer realidad el pilar europeo de derechos sociales es un compromiso político y una responsabilidad compartidos de las instituciones de la UE, las autoridades nacionales, regionales y locales, los interlocutores sociales y la sociedad civil. Porque una Europa social fuerte es la base no solo de la prosperidad y el bienestar de los ciudadanos de la UE, sino también de una economía competitiva.

Los dirigentes reafirmaron su compromiso con estos objetivos en la Cumbre Social de Oporto de mayo de 2021. Diferentes instituciones y organizaciones, incluido el Parlamento, adoptaron el Compromiso Social de Oporto, mientras que los jefes de Estado o de Gobierno de la Unión adoptaron la Declaración de Oporto. Los objetivos nacionales del pilar europeo de derechos sociales se presentaron en junio de 2022. Su aplicación se verá facilitada por el cuadro de indicadores sociales revisado y la financiación a través del plan de gasto de la Unión para el período 2021-2027 y NextGenerationEU, en particular el Mecanismo de Recuperación y Resiliencia. Se realizará un seguimiento en el marco del Semestre Europeo.

Entre las instituciones consultivas de la UE, y para el tema que se estudia, es importante mencionar al Comité Económico y Social. Este Comité es un órgano de consulta económica y social del Consejo de la Unión Europea y la Comisión Europea. Está compuesto por representantes de las organizaciones empresariales, de trabajadores y otros sectores representativos de la sociedad civil, especialmente de ámbitos socioeconómico, cívico, profesional y cultural.

\paragraph{Comité de Protección Social}
Por otro lado, el Comité de Protección Social. Se trata de un comité consultivo del Consejo de Empleo y Asuntos Sociales. Su objetivo es promover la cooperación entre las políticas de protección social de los Estados miembros y la Comisión. El CPS está compuesto por dos delegados de cada Estado miembro y de la Comisión.

El CPS se reúne todos los meses excepto en julio y agosto. Dos veces al año, las reuniones tienen lugar en el país que ejerce la Presidencia del Consejo de la UE.

El CPS elige a un presidente que ocupa el cargo por un período de dos años (renovable una vez). La presidenta actual es Ulrike Neufang (AT).

A su vez, el presidente está asistido por un , de laMesa, compuesto por la Comisión, cuatro vicepresidentes (dos vicepresidentes elegidos, un representante del país actual y un representante del siguiente país de la Presidencia del Consejo), el secretario del CPS, la Secretaría General del Consejo y el presidente del subgrupo sobre indicadores. El actual vicepresidente electo es Jere Päivinen (FI) y Marie Ranty (BE).

A través de un equipo de funcionarios, incluidos Katalin Szatmari , Paul Mintyy Ferdinand Egyed ., la Comisión proporciona laSecretaría del CPS.

El Comité de Protección Social (CPS) es un comité consultivo de política para los ministros del Consejo de Empleo y Asuntos Sociales (EPSCO). A raíz de una decisión del Consejo adoptada en virtud del artículo 160 del Tratado de Funcionamiento del. de la UE, se creó el CCP en 2004. La última Decisión del Consejo está en vigor desde 2015.

El Comité constituye un foro para la coordinación de políticas multilaterales, el diálogo y la cooperación a escala de la UE. Las principales funciones del CCP se definen como sigue (artículo 160del TFUE): 

supervisar la situación social y la evolución de las políticas de protección social en los Estados miembros y en la Unión;
promover el intercambio de información, experiencias y buenas prácticas entre los Estados miembros y con la Comisión;
elaborar informes, formular dictámenes o emprender otros trabajos en sus ámbitos de competencia, bien a petición del Consejo o de la Comisión, bien por propia iniciativa;
De conformidad con el artículo 148 del TFUE,, el Comité de Protección Social desempeña un papel clave en el contexto del Semestre Europeo y coopera estrechamente con otros órganos preparatorios del Consejo (especialmente el Comité de Empleo y el Comité de Política Económica) para debatir cuestiones de interés común, preparar los debates del Consejo sobre protección social y sobre las recomendaciones específicas por país.

Los ámbitos de competencia del CPS se reflejan en el marco del método abierto de coordinación social (MAC) y abarcan todos los principales capítulos de política social:

inclusión social
protección social
pensiones
cuidados de larga duración
asistencia sanitaria

De conformidad con su mandato basado en el Tratado, el CPS supervisa la situación social en la UE, en particular a través de su principal documento final, el Informe Anual del CPS. El informe proporciona al Consejo información sobre las principales prioridades de política social que deben tenerse en cuenta en el contexto de la preparación del próximo ciclo del Semestre Europeo. El informe:

hace un seguimiento de la situación social, incluidos los avances hacia el objetivo de reducción de la pobreza y la exclusión social para 2030, y destaca las tendencias más recientes que deben observarse;
identifica los principales retos sociales estructurales a los que se enfrentan los distintos Estados miembros, así como los buenos resultados sociales;
examina la evolución reciente de la política social;


\paragraph{Comité del Empleo}
El Comité de Empleo (EMCO) se crea en el año 2000, mediante una Decisión del Consejo basada en el artículo 150 del TFUE. Operaa en el marco político de la Estrategia Europea de Empleo y, tal como se establece en el artículo 148 del TFUE, desempeña un papel clave en el Semestre Europeo. Está compuesto por dos miembros de cada Estado miembro y de la Comisión y elige a un presidente que ocupa el cargo por un período de dos años (renovable una vez).\footnote{%
    El presidente es asesorado por un grupo directivo. Está compuesto por cuatro vicepresidentes (los presidentes de los subgrupos del Comité de Empleo y representantes de la Presidencia actual y entrante del Consejo de la UE), la Comisión, el secretario del Comité y la Secretaría General del Consejo de la UE.

La Secretaría del Comité de Empleo es proporcionada por la Comisión a través de un equipo de funcionarios que apoyan el trabajo del comité y organizan reuniones del Comité de Empleo.
}

Se trata del comité consultivo de los ministros de Empleo y Asuntos Sociales en el Consejo de Empleo y Asuntos Sociales (EPSCO) para promover la coordinación de las políticas de empleo y del mercado laboral tanto a nivel europeo como nacional. A excepción de julio y agosto, se reúne todos los meses del año y mantiene debates e intercambios políticos periódicos con otras partes interesadas clave, como los interlocutores sociales europeos y nacionales (sindicatos y empleadores) y la red de SPE y organizaciones como la OCDE y Eurofound. 

Cada año adopta su programa de trabajo teniendo en cuenta las prioridades políticas del Consejo Europeo, el Consejo y la Comisión. El programa de trabajo abarca cuatro ámbitos principales: (i) el Semestre Europeo (incluida la supervisión multilateral); (ii) trabajo temático; (iii) benchmarking; y, (iv) seguimiento/informes.

El trabajo del Comité de Empleo, en particular en lo que se refiere a la evaluación comparativa y el seguimiento/informes, cuenta con el apoyo del grupo de indicadores del Comité de Empleo, que tiene su propio programa de trabajo vinculado a las principales prioridades del Comité de Empleo y del Grupo de Evaluación del Comité de Empleo y aporta información a sus actividades.

\subsubsection{Dimensión Institucional}
Entre las instituciones fundamentales de la UE implicadas en asuntos de empleo y política social figuran el Consejo Europeo y la Comisión Europea. El primero determina las prioridades políticas de la UE. La segunda, entre otras tareas, gestiona las políticas de la UE de manera que presenta iniciativas conforme a los grandes objetivos establecidos por el Consejo Europeo. Además, la Comisión Europea se organiza en agencias y direcciones generales. Entre estas últimas figura la Direccion General de Empleo, Asuntos Sociales e Inclusión. Precisamente esta Dirección puso en marcha el Observatorio Europeo de las Políticas de Empleo. Este Observatorio funciona como un medio para facilitar infonnación sobre las políticas de empleo y las tendencias de los mercados de trabajo de la Unión. Además, también hay agencias de la UE enfocadas al tema que se estudia. Es el caso de la Fundación Europea para la Mejora de las Condiciones de Vida y de Trabajo. Se trata de una agencia que proporciona información, orientación y promoción para un entorno laboral seguro y sano.

\paragraph{Autoridad Laboral Europea}
Finalmente, se debe hacer mención a la La Autoridad Laboral Europea (en inglés, \href{https://www.hablamosdeeuropa.es/es/Paginas/Noticias/Autoridad-Laboral-Europea-(ELA).aspx}{European Labour Authority, ELA}) es el organismo de la Unión Europea responsable de coordinar y apoyar la aplicación del Derecho laboral europeo, garantizando que las normas de la UE sobre movilidad laboral y coordinación de la seguridad social se apliquen de manera justa y efectiva. La Oficina Europea de Coordinación de EURES (EC::O) está integrada en ELA. Fue creada en julio de 20 1 9 y desde septiembre de 2021 tiene su sede en Bratislava (Eslovaquia).

Par-a cumplir su objetivo supervisa las condiciones sociales en la Unión y el desarrollo de políticas de protección social en los países miembros. Además, los Estados miembros de la UE coordinan sus políticas teniendo en cuenta que todo ciudadano debe tener acceso a determinadas prestaciones sociales de manera que: a pesar de los diferentes regímenes de seguridad social, se proporcione una red de seguridad, se garanticen pensiones y atención sanitaria de calidad a un coste sostenible y se fomenté la inclusión social y la lucha contra la pobreza.

En 2019 se creó la Autoridad Laboral Europea (ALE), con sede en Bratislava, con el objetivo de ayudar a los Estados miembros y a la Comisión a velar por que las normas de la Unión sobre movilidad laboral y coordinación de la seguridad social se apliquen de manera justa, sencilla y eficaz. En 2026, la Comisión Europea presentará una propuesta para reforzar el papel de la ALE.


Posteriormente sustituidos por un Reglamento relativo a la aplicación de los regímenes de seguridad social a los trabajadores por cuenta ajena y a los miembros de sus familias que se desplazan dentro de la Comunidad (Reglamento (CEE) n.o 1408/71). Los ciudadanos de Islandia, Liechtenstein y Noruega están cubiertos por el Acuerdo del Espacio Económico Europeo (EEE) y los de Suiza por el Acuerdo UE-Suiza. Tras el Brexit, se siguen protegiendo los derechos de las personas cubiertas por el Acuerdo de Retirada celebrado entre la Unión Europea y el Reino Unido. En el caso de las personas no cubiertas por el Acuerdo de Retirada, la coordinación de la seguridad social se regula mediante el correspondiente Protocolo del Acuerdo de Comercio y Cooperación.

\paragraph{Centro Europeo para el Desarrollo de la Formación Profesional (CEDEFOP, 1975)}
El CEDEFOP, fundada en 1975 y con sede en Grecia desde 1995, apoya la promoción, el desarrollo y la aplicación de políticas de EFP, capacidades y cualificaciones en la Unión Europea. Mejora y difunde conocimientos, proporciona datos y servicios basados en la investigación para la elaboración de políticas y facilita el intercambio de conocimientos entre los agentes nacionales y de la Unión. Ha desarrollado enfoques a escala de la Unión para comprender la oferta y la demanda de capacidades actuales y anticiparse al futuro.

La Unión apoya la EFP principalmente con financiación en el marco del Fondo Social Europeo Plus (FSE+) (3.6.2) y Erasmus+. En los primeros años, hasta 1970, el FSE reembolsaba a los Estados miembros el $50\%$ de los costes de formación profesional y de las indemnizaciones de reinstalación de los trabajadores víctimas de reestructuraciones económicas. Prestó asistencia a más de dos millones de personas durante ese período. Desde entonces se han revisado sus funciones y normas de funcionamiento y se han integrado instrumentos complementarios de la Unión. En consecuencia, el FSE+ se ha convertido en el principal instrumento de la Unión para invertir en las personas. El presupuesto total del FSE+ en el actual período de programación 2021-2027 asciende a casi 99 300 millones EUR. Los objetivos específicos del FSE+ son los siguientes: — mejorar la calidad, la inclusividad, la eficacia y la pertinencia de los sistemas de educación y de formación para el mercado de trabajo, también mediante la validación del aprendizaje no formal e informal; — promover la igualdad de acceso a una educación y formación inclusivas y de calidad, así como la igualdad a la hora de terminar los estudios correspondientes, en particular en el caso de los grupos desfavorecidos, facilitando la movilidad educativa para todos y la accesibilidad para las personas con discapacidad, entre otras cosas;
— promover el aprendizaje permanente, en particular mediante oportunidades para todos de perfeccionamiento y reciclaje flexibles de las capacidades teniendo en cuenta las capacidades empresariales y digitales, una mejor previsión de los cambios y nuevos requisitos de capacidades, habida cuenta de las necesidades del mercado de trabajo, facilitando las transiciones profesionales y promoviendo la movilidad profesional.

Los Estados miembros con una elevada proporción media de jóvenes (de entre 15 y 29 años) que ni estudian, ni trabajan, ni reciben formación (ninis) están obligados a gastar el 12,5 % de sus recursos del FSE+ en régimen de gestión compartida en apoyo del empleo juvenil, la EFP, la transición de la escuela al trabajo, los itinerarios para reintegrarse en la educación o la formación y la educación de segunda oportunidad, en particular en el contexto de la ejecución de los programas en el marco de la Garantía Juvenil.

El programa Erasmus+ cuenta con un presupuesto de 26 600 millones EUR para el período 2021-2027 (frente a los 14 700 millones EUR del período anterior). A través del aprendizaje permanente, apoya el desarrollo educativo, profesional y personal de las personas en los ámbitos de la educación y la formación, la juventud y el deporte, y promueve la movilidad educativa. La edición actual busca ampliar su alcance para incluir a las personas de entornos económicos desfavorecidos, con discapacidad o dificultades educativas, así como a personas que viven en zonas rurales y remotas. Cada año, más de 130 000 estudiantes y aprendices de EFP, así como 20 000 miembros del personal de EFP, se benefician de oportunidades de movilidad Erasmus+. El programa también financia centros de excelencia profesional, alianzas para la innovación, asociaciones de cooperación, proyectos de desarrollo de capacidades con países no pertenecientes a la Unión y proyectos prospectivos en el ámbito de la EFP. En las orientaciones políticas de la Comisión para el período 2024-2029, la presidenta de la Comisión se comprometió a reforzar Erasmus+, también para la formación profesional, y a aumentar el número de personas que se benefician del programa.



\clearpage
\section{Conclusión} 


\subsection{Recapitulación}


\subsection{Extensiones y relación con otras partes del Temario}



\subsection{Opinión}

\subsubsection{Idea final (Salida o cierra)}

\clearpage
\pagenumbering{gobble}
\MyBibliography



\href{https://www.cepc.gob.es/sites/default/files/2021-12/35219rei07002067.pdf}{\textit{La política regional de la CEE. Fundamemos e instrumentación.} PALAU (1986). Revista de Estudios Internacionales. Vol. 7, Nº2}

La comunitarización de la política regional. SUÁREZ et al. (1996), Boletín económico de ICE. Nº 2519. 1 996


\href{https://www.revistarue.eu/RUE/092007.pdf}{Cincuenta años en la configuración de la Política Regional Europea: 1957-2007. SUÁREZ et. al. (2007). Revista Universitaria Europea, Nº9}

\href{http://www.revista-eea.net/documentos/27312.pdf}{\textit{El desarrollo histórico de la Política Social de la Unión Europea y su estado actual}. SANZ (2009). Estudios de Economía Aplicada. Vol. 27-3}


\MyBibItem{DAWID y GATTI}{Chapter 2 - Agent-Based Macroeconomics}{2018}{https://doi.org/10.1016/bs.hescom.2018.02.006}




% ANEXOS
\clearpage
\anexo{\textbf{Preguntas Test}}
%\input{\TemaCode/questions.tex} 


\clearpage
\anexo{Implementación: Pilar Europeo de Derechos Sociales}\label{anx:consolidacion_peds}
El pilar social tiene por objeto impulsar una Europea social para toda la ciudadanía europea. Aspira a reforzar el acervo social y a hacer más efectivos los derechos de los ciudadanos. Presta especial atención a los aspectos sociales y de empleo, y a velar por que el modelo social europeo esté preparado para afrontar los desafíos del siglo XXI. El objetivo del pilar es contribuir al progreso social promoviendo mercados de trabajo y sistemas de protección social justos y que funcionen correctamente.
    
En la Cumbre Social de Gotemburgo de 2017, el Parlamento, el Consejo y la Comisión pusieron de relieve su compromiso compartido mediante la adopción de unaproclamación conjunta sobre el pilar europeo de derechos sociales. El pilar europeo de derechos sociales va acompañado de un «cuadro de indicadores sociales» para el seguimiento de los avances.
    
En la Cumbre Social de Gotemburgo de 2017, el Parlamento, el Consejo y la Comisión pusieron de relieve su compromiso compartido mediante la adopción de unaproclamación conjunta sobre el pilar europeo de derechos sociales. El pilar europeo de derechos sociales va acompañado de un «cuadro de indicadores sociales» para el seguimiento de los avances.



\clearpage
\anexo{Política Social y de Empleo de la Unión Europea: Dimension jurídica}\label{anx:normativa_empleo}

\textbf{Consultar:} \href{https://www.insst.es/documents/94886/3978182/Tema+12.+La+Normativa+de+la+UE+en+Materia+Social+y+de+Empleo.pdf/a17600df-3f07-6bb4-68c4-712b0fc6071f?t=1663571188948}{Tema 12. Primer Ejercicio del Proceso Selectivo a INSST (2023)}

La normativa de la Unión (UE) en materia social y de empleo tiene su origen en los Tratados cuyos principios y objetivos en esta materia han producido un completo corpus legislativo, derecho derivado. Así, la dimensión social de la UE se ha ido desarrollando gradualmente a lo largo del proceso de integración europea.

La UE sólo tiene competencias limitadas en el ámbito social, por lo que la responsabilidad del empleo y las políticas sociales recae principalmente en los gobiernos nacionales. En consecuencia, la UE ha creado una serie de instrumentos en el sector social que incluyen leyes, fondos o herramientas para coordinar y supervisar mejor las políticas nacionales. No obstante, la Unión toma medidas para garantizar la coordinación de las políticas de empleo de los Estados miembros, en particular definiendo las orientaciones de dichas políticas y alentando a los Estados miembros a compartir las mejores prácticas en sectores como la inclusión social, la pobreza y las pensiones.

\textbf{Empleo}. Los principales aspectos en materia de empleo que define el Tratado se recogen en el título IX (artículos 145 a 150). Los Estados miembros y la Unión se esforzarán por desarrollar una estrategia coordinada para el empleo, en particular para potenciar una mano de obra cualificada, formada y adaptable y lograr mercados laborales con capacidad de respuesta al cambio económico. Así, los Estados miembros desarrollarán sus políticas de empleo de forma compatible con las orientaciones generales de las políticas económicas de los Estados miembros y de la Unión, considerarán el fomento del empleo como un asunto de interés común y coordinarán sus actuaciones al respecto en el seno del Consejo. Para la consecución del objetivo de un alto nivel de empleo, la Unión promueve la cooperación entre los Estados miembros, apoya y, completa sus respectivas actuaciones respetando las competencias de los Estados miembros. Anualmente el Consejo Europeo examina la situación del empleo en la Unión y adopta conclusiones al respecto, basándose en un informe conjunto anual elaborado por el Consejo y la Comisión. De esta forma se elaboran las orientaciones que los Estados miembros tendrán en cuenta en sus respectivas políticas de empleo. Por otra parte, cada Estado miembro facilita al Consejo y a la Comisión un informe anual sobre las principales medidas adoptadas para aplicar su política de empleo, a la vista de las orientaciones de empleo. El Consejo efectúa un examen anual de la aplicación de las políticas de empleo de los Estados miembros y de esta forma formula recomendaciones a los Estados miembros, si lo considera pertinente a la vista de dicho examen. También con carácter anual el Consejo y la Comisión elaboran un informe conjunto para el Consejo Europeo sobre la situación del empleo en la Unión y sobre la aplicación de las orientaciones para el empleo. El Parlamento Europeo y el Consejo, con arreglo al procedimiento legislativo ordinario y previa consulta al Comité Económico y Social y al Comité de las Regiones, podrán adoptar medidas para alentar la cooperación entre los Estados miembros y apoyar la actuación de estos últimos en el ámbito del empleo, a través de iniciativas destinadas a desarrollar los intercambios de información y buenas prácticas, facilitar análisis comparativos y asesoramiento, así como promover planteamientos innovadores y evaluar experiencias, en particular recurriendo a proyectos piloto. Estas medidas no incluirán armonización alguna de las disposiciones legales y reglamentarias de los Estados miembros. El Consejo, por mayoría simple, previa consulta al Parlamento Europeo, creará un Comité de Empleo de carácter consultivo para fomentar la coordinación entre los Estados miembros en materia de políticas de empleo y del mercado laboral. Este Comité supervisa la situación del empleo y las políticas en materia de empleo de los Estados miembros y de la Unión y, previa consulta a los interlocutores sociales, elabora dictámenes y contribuye a la preparación de las medidas del Consejo anteriormente citadas.


Las normas de derecho derivado en materia de trabajo son (\textcolor{yellos}{10/07/2023}):
\begin{lnum}
    \item Directiva del Consejo, de 12 de junio de 1989, relativa a la aplicación de medidas para promover la mejora de la seguridad y de la salud de los trabajadores en el trabajo (transpuesta por la Ley de Prevención de Riesgos Laborales)
    \item Directiva 89/656/CEE del Consejo, de 30 de noviembre de 1989, relativa a las disposiciones mínimas de seguridad y de salud para la utilización por los trabajadores en el trabajo de equipos de protección individual
    \item Directiva 94/33/CE, sobre protección de los jóvenes en el trabajo establece las disposiciones mínimas a fin de garantizar un mayor nivel de protección de la seguridad y de la salud de las personas trabajadoras jóvenes (transpuesta por la Ley de Prevención de Riesgos Laborales)
    \item Directiva 1999/92/CE del Parlamento Europeo y del Consejo, de 16 de diciembre de 1999, relativa a las disposiciones mínimas para la mejora de la protección de la salud y la seguridad de los trabajadores expuestos a los riesgos derivados de atmósferas explosivas (transpuesta por el Real Decreto 681/2003)
    \item Directiva 2009/104/CE del Parlamento Europeo y del Consejo, de 16 de septiembre de 2009, relativa a las disposiciones mínimas de seguridad y de salud para la utilización por los trabajadores en el trabajo de los equipos de trabajo (transpuesta por la Ley de Prevención de Riesgos Laborales)
    \item Reglamento (UE) 2016/425 del Parlamento Europeo y del Consejo, de 9 de marzo de 2016, relativo a los equipos de protección individual y por el que se deroga la Directiva 89/686/CEE del Consejo.
    \item Directiva (UE) 2019/1152 del Parlamento Europeo y del Consejo, de 20 de junio de 2019, relativa a unas condiciones laborales transparentes y previsibles en la Unión Europea (aún pendiente de transposición)
    \item Directiva 79/7/CEE relativa a la aplicación progresiva del principio de igualdad de trato entre hombres y mujeres en materia de seguridad social (aplicación directa)
    \item Directiva 2000/78/CE relativa al establecimiento de un marco general para la igualdad de trato en el empleo y la ocupación. (transpuesta por la Ley 51/2003 y Ley 62/2003).
    \item Directiva 2006/54/CE relativa a la aplicación del principio de igualdad de oportunidades e igualdad de trato entre hombres y mujeres en asuntos de empleo y ocupación. (transpuesta por la Ley Orgánica 3/2007)
    \item Directiva 2003/88/CE del Parlamento Europeo y del Consejo, de 4 de noviembre de 2003, relativa a determinados aspectos de la ordenación del tiempo de trabajo (transpuesta por el Real Decreto 311/2016).
    \item Directiva (UE) 2019/1158 del Parlamento Europeo y del Consejo, de 20 de junio de 2019, relativa a la conciliación de la vida familiar y la vida profesional de los progenitores y los cuidadores (transpuesta por Real Decreto-ley 6/2019 y Real Decreto-Ley 5/2023)
    \item Directiva (UE) 2022/2041 del Parlamento Europeo y del Consejo de 19 de octubre de 2022 sobre unos salarios mínimos adecuados en la Unión Europea (todavía no transpuesta)
    \item Directiva 98/59/CE del Consejo, de 20 de julio de 1998, relativa a la aproximación de las legislaciones de los Estados miembros que se refieren a los despidos colectivos (transpuesta por el Estatuto de los Trabajadores, art. 51)
    \item Directiva 2008/104/CE del Parlamento Europeo y del Consejo, de 19 de noviembre de 2008, relativa al trabajo a través de empresas de trabajo temporal (transpuesta por la Ley 35/2010)
    \item Directivas 92/56/CEE y 98/59/CE del Consejo, por la que se exige que, en caso de despidos masivos, los empresarios entablen negociaciones con los trabajadores (transpuesta por el Estatuto de los Trabajadores, art.51 y Real Decreto 1483/2012)
    \item Directiva 94/45/CE del Consejo, modificada por la Directiva 2009/38/CE, sobre la constitución de comités de empresa europeos. La legislación sobre estos comités de empresa abarca las empresas multinacionales con al menos 1 000 trabajadores en la Unión o el Espacio Económico Europeo y por lo menos 150 trabajadores en dos o más Estados miembros. (transpuesta por la Ley 44/1999)
    \item Directiva 2001/23/CE del Consejo relativa al mantenimiento de los derechos de los trabajadores en caso de traspasos de empresas, por la que se estipula que los trabajadores deben ser informados de las razones del traspaso y de sus consecuencias (transpuesta por el artículo 44 del Estatuto de los Trabajadores)
    \item Directiva 2001/86/CE del Consejo por la que se completa el Estatuto de la Sociedad Anónima Europea en lo que respecta a la implicación de los trabajadores y se fijan normas para la participación de estos en las decisiones sobre el desarrollo estratégico de la empresa (transpuesta por la Ley 31/2006)
    \item Directiva 2002/14/CE por la que se establece un marco general relativo a la información y a la consulta de los trabajadores en la Comunidad Europea, que fija normas mínimas de procedimiento (transpuesta por la Ley 38/2007)
    \item Directiva 2003/72/CE del Consejo por la que se completa el Estatuto de la sociedad cooperativa europea en lo que respecta a la implicación de los trabajadores y se garantiza que los representantes de estos puedan influir en la gestión de las sociedades cooperativas europeas (transpuesta por la Ley 31/2006)
    \item Directiva 2004/25/CE relativa a las ofertas públicas de adquisición, en virtud de la cual debe preverse la posibilidad de que los trabajadores de las sociedades afectadas o sus representantes puedan pronunciarse sobre las posibles repercusiones de la oferta en cuestión sobre el empleo. También son de aplicación las normas generales en materia de información y consulta de los trabajadores (transpuesta por la Ley 6/2007 y Real Decreto 1066/2007)
    \item Directiva (UE) 2019/2121 en lo que atañe a las transformaciones, fusiones y escisiones transfronterizas. La Directiva modificada refuerza la posición de los trabajadores y sus representantes en materia de información, consulta y derechos de participación antes de una transformación transfronteriza. Impone a la dirección la obligación de responder a los sindicatos y a las organizaciones de trabajadores y consultarles sobre las consecuencias de la transformación prevista. (transpuesta por el Real Decreto-ley 5/2023)
    \item Directiva (UE) 2024/2831 del Parlamento Europeo y del Consejo, de 23 de octubre de 2024, relativa a la mejora de las condiciones laborales en el trabajo en plataformas;
    \item \textbf{Reglamentos de Coordinación de Seguridad Social} (Reglamento 883/2004): Protegen los derechos de salud, pensiones y desempleo de los ciudadanos que se desplazan dentro de la UE, el Espacio Económico Europeo y Suiza. Para concoer más sobre la materia, se puede consultar esta publicación: \href{https://ec.europa.eu/employment_social/soc-prot/schemes/guide_es.pdf}{\textit{Normativa Comunitaria en materia de Seguridad Social}. Comisión Europea (1999)}
    \item Directiva de Trabajadores Desplazados (Directiva 96/71/CE y modificada por la Directiva 2018/957): Establece que los empleados enviados temporalmente a otro país miembro deben percibir el salario mínimo del país de acogida y cumplir con sus normativas de jornada laboral y descansos
    \item Además, las directivas relativas a seguridad laboral y salud en el trabajo, labor esencial de la Agencia Europea para la Seguridad y Salud en el Trabajo: \href{https://osha.europa.eu/en/safety-and-health-legislation/european-directives}{European directives on safety and health at work} 
\end{lnum}




\textbf{Política Social}. Los principales aspectos en materia de política social que define el Tratado están recogidos,
principalmente en el título X (artículos 151 a 161). La Unión y los Estados miembros, tendrán como objetivo el fomento del empleo, la mejora de las condiciones de vida y de trabajo, a fin de conseguir su equiparación por la vía del progreso, una protección social adecuada, el diálogo social, el desarrollo de los recursos humanos para conseguir un nivel de empleo elevado y duradero y la lucha contra las exclusiones. Para ello, emprenderán acciones en las que se tenga en cuenta la diversidad de las prácticas nacionales, en particular en el ámbito de las relaciones contractuales, así como la necesidad de mantener la competitividad de la economía de la Unión.La Unión reconoce y promueve el papel de los interlocutores sociales en su ámbito y facilita el diálogo entre ellos. La cumbre social tripartita para el crecimiento y el empleo contribuye al diálogo social. La Unión apoya y completa la acción de los Estados miembros en los siguientes ámbitos: la mejora, en concreto, del entorno de trabajo, para proteger la salud y la seguridad de las personas trabajadoras; las condiciones de trabajo; la seguridad social y la protección social de las personas trabajadoras; su protección en caso de rescisión del contrato de trabajo; la información y la consulta; la representación y la defensa colectiva de los intereses de la población trabajadora y del empresariado; las condiciones de empleo de los nacionales de terceros países que residan legalmente en el territorio de la Unión; la integración de las personas excluidas del mercado laboral; la igualdad entre hombres y mujeres por lo que respecta a las oportunidades en el mercado laboral y al trato en el trabajo; la lucha contra la exclusión social; la modernización de los sistemas de protección social. A tal fin, el Parlamento Europeo y el Consejo adoptan medidas destinadas a fomentar la cooperación entre los Estados miembros mediante iniciativas para mejorar los conocimientos, desarrollar el intercambio de información y de buenas prácticas, promover fórmulas innovadoras y evaluar experiencias, con exclusión de toda armonización de las disposiciones legales y reglamentarias de los Estados miembros. Asimismo, mediante directivas adoptan en determinados ámbitos las disposiciones mínimas de aplicación en cada uno de los Estados miembros. El Parlamento Europeo y el Consejo decidirán, en función del ámbito correspondiente, con arreglo al procedimiento legislativo ordinario o especial.

Las normas de derecho derivado en materia social son (\textcolor{yellos}{10/07/2023}):
\begin{lnum}
    \item Directiva 2000/43/CE relativa a la aplicación del principio de igualdad de trato de las personas independientemente de su origen racial o étnico (transpuesta por la Ley 62/2003)
    \item Directiva 2004/113/CE por la que se aplica el principio de igualdad de trato entre hombres y mujeres al acceso a bienes y servicios y su suministro (transpuesta por la Ley Orgánica 3/2007).
\end{lnum}





La Recomendación del Consejo sobre el desarrollo de condiciones marco para la economía social se adoptó en noviembre de 2023 para avanzar hacia los objetivos del pilar europeo de derechos sociales, ayudando a los Estados miembros a adaptar su legislación y sus políticas a las necesidades de una economía social.

Dado que las plataformas digitales de trabajo emplean a una proporción cada vez mayor de trabajadores, a menudo con horarios irregulares y sometidos a gestión algorítmica, el Parlamento y el Consejo han adoptado la Directiva relativa a la mejora de las condiciones laborales en el trabajo en plataformas digitales (Directiva (UE) 2024/2831, de 23 de octubre de 2024).

\clearpage
\anexo{Agencias Europeas: En materia de Cohesión, Social y de Empleo}\label{anx:organismosue_3b40}

A continuación se listan las organizaciones vinculadas, directa o indirectamente, a las Instituciones Europeas o sus actividades, cuyo ámbito de actuación se encuadra en el objeto de este Tema:

\href{https://ec.europa.eu/info/departments/european-research-executive-agency_es}{\textbf{Agencia Ejecutiva Europea de Investigación (AEI)}}. La Agencia Ejecutiva de Investigación (AEI) es un organismo de financiación para investigación e innovación. Gestiona las subvenciones de la UE para investigación.

\href{https://osha.europa.eu/es}{\textbf{Agencia Europea para la Seguridad y la Salud en el Trabajo (EU-OSHA)}}. La Agencia Europea para la Seguridad y la Salud en el Trabajo contribuye a que los lugares de trabajo sean más seguros, saludables y productivos.

\href{https://www.ela.europa.eu/es}{\textbf{Autoridad Laboral Europea}} (Agencia descentralizada). La Autoridad Laboral Europea vela por que el trabajo transfronterizo y las normas de seguridad social se apliquen correctamente en la UE.

\href{https://www.eib.org/en/index.htm}{\textbf{Banco Europeo de Inversiones}} (Órgano de la UE). El Banco Europeo de Inversiones, en calidad de banco para el clima de la UE, ofrece financiación y asesoramiento especializado para proyectos cuyo objetivo sea luchar contra el cambio climático, fomentar la integración europea, promover el desarrollo de la UE y apoyar las políticas de la UE en más de 160 países.

\href{https://www.cedefop.europa.eu/es}{\textbf{Centro Europeo para el Desarrollo de la Formación Profesional (Cedefop)}} (Agencia descentralizada). El Cedefop ayuda a las instituciones y organizaciones de la UE a desarrollar y satisfacer las necesidades de formación profesional.

\href{https://cinea.ec.europa.eu}{\textbf{CINEA}} (Agencia ejecutiva). La Agencia Ejecutiva Europea de Clima, Infraestructuras y Medio Ambiente (CINEA) gestiona los programas de la Comisión Europea que contribuyen a la descarbonización y al crecimiento sostenible.

\href{https://cor.europa.eu/es/Pages/default.aspx}{Comité Europeo de las Regiones (CDR)} (Órgano de la UE). El Comité Europeo de las Regiones es una asamblea que asesora a las instituciones europeas sobre las leyes que afectan a las ciudades y regiones de la UE.

\href{https://commission.europa.eu/about-european-commission/departments-and-executive-agencies/european-education-and-culture-executive-agency_es}{EACEA} (Agencia ejecutiva). La Agencia Ejecutiva Europea de Educación y Cultura (EACEA) gestiona la financiación en educación, cultura, sector audiovisual, deporte, ciudadanía y voluntariado.

\href{https://ec.europa.eu/info/departments/european-research-council_es}{\textbf{ERCEA}} (Agencia ejecutiva). La Agencia Ejecutiva del Consejo Europeo de Investigación gestiona los programas de la UE en materia de investigación.

\href{https://www.etf.europa.eu/es}{\textbf{Fundación Europea de Formación (ETF)}} (Agencia descentralizada). La Fundación Europea de Formación apoya a los países socios de la UE en la reforma de sus sistemas de educación y formación.

\href{https://ec.europa.eu/info/departments/european-health-and-digital-executive-agency_es}{\textbf{HADEA}} (Agencia ejecutiva). La Agencia Ejecutiva Europea en los ámbitos de la Salud y Digital (HaDEA) encarna la ambición de la Comisión Europea de contribuir a reconstruir Europa tras la COVID-19 desde su núcleo, de manera que sea más ecológica, más digital, más resiliente y mejor adaptada a los retos actuales y futuros.

\href{https://eit.europa.eu/}{\textbf{Instituto Europeo de Innovación y Tecnología (EIT)}} (Agencia descentralizada). El Instituto Europeo de Innovación y Tecnología tiene por objeto aumentar la innovación mediante la cooperación entre la industria y las organizaciones educativas.

\href{https://eige.europa.eu/}{\textbf{Instituto Europeo de la Igualdad de Género (EIGE)}} (Agencia descentralizada). El Instituto Europeo de la Igualdad de Género contribuye a los esfuerzos por lograr la igualdad de género en la UE.




\end{document}